%%%%%%%%%%%%%%%%%%%%%%%%%%%%%%%%%%%%%%%%%%%%%%%%%%%%%%%%%%
% # L2 Research on Relative Clauses and L1/L2 Pathways for Learning
%%%%%%%%%%%%%%%%%%%%%%%%%%%%%%%%%%%%%%%%%%%%%%%%%%%%%%%%%%

\chapter{L2 Research on Relative Clauses and L1/L2 Pathways for Learning}\label{ch-acquisition}

% \excnt=102

This chapter introduces the reader to previous research on L2 resumption in RCs and explores how learners develop L1 vs. L2 implicit knowledge about the distribution of gapped vs. resumptive RCs. It begins in Sections~\ref{acq-findings}--\ref{acq-hawkins} by reviewing some key findings from the relevant L2 experimental literature. It then proceeds to identify several areas that are underexplored in the existing research in Sections~\ref{acq-slw}--\ref{acq-assumption} and to consider how the acquisition of RCs might unfold in L1 vs. interlanguage in Sections~\ref{acq-preliminary}--\ref{acq-initial-state}. Along the way, I~introduce several of the central goals of the dissertation. Section~\ref{acq-conclusion} concludes the chapter.

%========================================================%
\section{Key findings from the L2 experimental literature on RCs}\label{acq-findings}
%========================================================%

Previous experimental research on L2 resumption in RCs has generally found that it is more frequently produced in environments thought to be difficult for relativization (\cite{algady2013}; \cite{eckman.etal1988}; \cite{gass1979}; \cite{hyltenstam1984,hyltenstam1990}; \cite{kim2013}; \cite{maghrabi1997}; \cite{pavesi1986}). Some studies have also shown that in cases where the target language lacks grammatical resumption, low-proficiency L2ers produce and accept resumptive RCs at elevated rates compared to high-proficiency L2ers (\cite{hawkins.chan1997}; \cite{kim2013}; \cite{maghrabi1997}). Studies showing that L2ers produce resumptive RCs irrespective of their well-formedness in the L1 and the target language (\cite{gass1979}; \cite{hyltenstam1984,hyltenstam1990}; \cite{maghrabi1997}; \cite{pavesi1986}) are of particular interest because they reveal that L2 resumption cannot always be traced in a straightforward way to either L1 transfer or target language input.

%========================================================%
\section{Gass (1979)}\label{acq-gass}
%========================================================%

Noticing crosslinguistic patterns in the grammatical functions that can be targeted by single-clause relativization, \textcite{keenan.comrie1977} formulated the Noun Phrase Accessibility Hierarchy (NPAH) in (102), according to which functions in higher positions (further to the left) are more widely available for relativization than those in lower positions (further to the right).

\pex 
SU > DO > IO > OBL > GEN > OCOMP\footnotemark
\xe 

\footnotetext{As a reminder: SU = subject; DO = direct object; IO = indirect object; OBL = oblique; GEN = genitive; OCOMP = object of comparison}

In the years following the publication of this Keenan and Comrie paper, there was a flurry of studies seeking to test whether the NPAH---initially conceived as typological generalizations---is also predictive of developmental patterns in the L2 acquisition of single-clause RCs (\cite{eckman.etal1988}; \cite{gass1979}; \cite{hawkins.chan1997}; \cite{ioup.kruse1977}; \cite{schachter.etal1976}; \cite{tarallo.myhill1983}; \cite{yip.matthews1991}). The focus of these studies, at least at first, was not on gapping vs. resumption \textit{per se} but rather on L2 accuracy in the production and acceptance of grammatical RCs at various positions on the hierarchy. 

\textcite{gass1979} was one of the first studies in this line of research. She administered an acceptability judgment task (AJT) and a sentence combining task (SCT) to adult L2ers of English from nine different L1 backgrounds (Arabic, Chinese, French, Italian, Japanese, Korean, Persian, Portuguese, and Thai), as detailed in Table \ref{tblr:gass}. Gass's primary aim was to test whether L2ers' accuracy in producing and judging RCs would decline down the NPAH.

\begin{longtblr}[caption={\textcite{gass1979}: Experimental design}, label={tblr:gass}]{colspec={|X|}, rowhead=0, rows={font=\raggedright\setstretch{1}}}\hline
\textbf{Participants:} 17 high-intermediate to advanced adult L2ers of English distributed across 9 L1s (Arabic, Chinese, French, Italian, Japanese, Korean, Persian, Portuguese, and Thai), the assumption being that only Arabic, Chinese, and Persian have grammatical resumption; it is unclear how many participants there were for each L1 \vspace{5pt}\rule{0pt}{10pt}\\*
\textbf{Tasks:} Each participant took part in an acceptability judgment task (AJT) and a sentence combining task (SCT); participants also completed a third task for a different study; the ordering of the tasks is not specified; each participant completed the tasks six times over the course of four months with different lexicalizations each time \vspace{5pt}\rule{0pt}{10pt}\\
\textbf{AJT stimuli:} 29 sentences with restrictive RCs (13 well-formed, 16 ill-formed); it is unclear exactly how many of the ill-formed sentences contained RPs, but there was at least one RP trial each for SU, DO, IO, GEN, and OCOMP (OBL appears to have been omitted)\vspace{5pt}\\
\textbf{AJT procedure:} Participants read sentences and assigned binary acceptability ratings to them \vspace{5pt}\\
\textbf{SCT stimuli:} 12 sentence pairs manipulating the positions of the italicized NP in the first sentence (SU or DO) and the identical NP in the second sentence (SU, DO, IO, OBL, GEN, or OCOMP) \vspace{5pt}\\
\textbf{SCT procedure:} Participants were given pairs of sentences and instructed to combine them via relativization \vspace{5pt}\\
\textbf{Example SCT item in the [SU, DO] condition (p. 333, (2)):} \vspace{5pt}\\*
\setlength{\leftskip}{10pt}\textbf{Stimulus:} \par\nobreak
\textit{The girl} ran home. I saw \textit{the girl}.
\par\nobreak\vspace{5pt} 
\textbf{Target response:} \par\nobreak
The girl [that I saw \_\_] ran home.
\\\hline
\end{longtblr}

Gass's participants produced and accepted resumptive RCs at least some of the time at all positions on the hierarchy, even in subject position (which, recall, is universally disallowed). Participants with L1s that have grammatical resumption (Arabic, Chinese, and Persian) tended to produce and accept resumptive RCs at higher rates, but this difference came out as statistically significant for only the top three positions on the hierarchy (SU, DO, and IO); presumably, the difference was not significant at the lower positions because participants in all groups produced and accepted resumptive RCs at fairly high rates in those environments. Gass took these results as evidence that Arabic, Chinese, and Persian speakers transferred resumption from their L1 to their interlanguage English.

These results broadly indicate that it is not uncommon for L2ers of English, regardless of L1 background, to produce and/or accept RPs in RCs. The results also suggest that resumption can be transferred from the L1 to the interlanguage. However, the study has a number of important limitations. For one thing, since each L1 group had perhaps only one or two participants, it would be hasty to make generalizations about any of the sampled populations. Second, because there was no control group of native English speakers, we cannot be certain what nativelike performance on the tasks in regard to resumption would have looked like. Third, the AJT was not specifically designed to examine resumption, and as such it had only a few trials containing resumptive RCs. Fourth, the SCT employed a rather contrived procedure for eliciting RCs (i.e.,~participants were presented with two sentences and asked to combine them), and this may have increased the likelihood of participants drawing on explicit grammatical knowledge when formulating responses. Finally, no proficiency measure was included in the experiment, so we have no way of knowing how participants compared in terms of English proficiency. In light of these problems, it is wise to interpret the results with caution.

%========================================================%
\section{Hyltenstam (1984, 1990)}\label{acq-hyltenstam}
%========================================================%

\textcite{hyltenstam1984,hyltenstam1990} was a pathbreaking study in this line of research as well as one of the primary inspirations for the present dissertation. Unlike the studies that came before it, Hyltenstam's was designed specifically to test rates of gapping vs. resumption in the L2\footnotemark{} production of single-clause RCs. It involved a Swedish RC elicitation task with conditions targeting the six positions on the NPAH, as shown in Table \ref{tblr:hyltenstam}. 

\footnotetext{The participants had likely learned at least some English prior to starting on Swedish; I follow \textcite{hyltenstam1984,hyltenstam1990} in using `L2' for nonnative languages in general. Thus, any nonnative learner can be described as an L2er.}

\begin{longtblr}[caption={\textcite{hyltenstam1984,hyltenstam1990}: Oral RC elicitation task}, label={tblr:hyltenstam}]{colspec={|X|}, rowhead=0, rows={font=\raggedright\setstretch{1}}}\hline
\textbf{Participants:} 45 advanced adult L2ers of Swedish distributed across four L1s (9~Finnish, 12~Greek, 12~Persian, and 12~Spanish), the assumption being that Greek and Persian have grammatical resumption whereas Swedish, Finnish, and Spanish do not \vspace{5pt}\rule{0pt}{10pt}\\*
\textbf{Stimuli:} There were six conditions eliciting relativization from different positions on the NPAH~(SU, DO, IO, OBL, GEN, and OCOMP); each participant received 36 items (6~per~condition) \vspace{5pt}\\
\textbf{Procedure:} The experimenter showed the participant drawings of numbered women, men, girls, and boys taking part in various events and asked the participant questions about them (e.g.,~\textit{Who is number five?}) \vspace{5pt}\\
\textbf{Example response with a resumptive RC in the GEN condition (\cite{hyltenstam1984}, p.~54):} \vspace{5pt}\\*
\setlength{\leftskip}{10pt}\begingl
\gla \ljudge*Det är mannen som hunden biter i hans väska. //
\glb it is the.man that the.dog is.biting in his bag //
\glft `It is the man whose bag the dog is biting.' //
\endgl
\\\hline
\end{longtblr}

The participants were advanced adult learners of Swedish distributed across four L1s: Finnish, Greek, Persian, and Spanish. These languages were selected because they differ regarding the availability of resumption as a licit option for relativization in single-clause RCs. Table \ref{tab:hyltenstam} shows that, according to \textcite{keenan.comrie1977}, Greek and Persian have grammatical resumption whereas Swedish, Finnish, and Spanish do not. Hyltenstam had no independent measure of Swedish proficiency, so he was not able to compare participant's performance in that regard.

\begin{table}
\caption{Distribution of gapping (−), resumption (+), and unrelativizable domains (0) in single-clause RCs for the languages in \textcite{hyltenstam1984,hyltenstam1990}}\label{tab:hyltenstam}
\raggedright\setstretch{1}
\begin{tabularx}{\linewidth}{lZZZZZZ}\hline
Language & SU & DO & IO & OBL & GEN & OCOMP \\\hline
Swedish & − & − & − & − & − & − \\
Finnish & − & − & − & − & − & 0 \\
Spanish & − & − & − & − & − & 0 \\
Greek & − & − & +? & +? & + & + \\
Persian & − & ± & + & + & + & + \\\hline
\end{tabularx}\par\nobreak\vspace{5pt}
\hfill (adapted from \cite{hyltenstam1984}, p.~46, Table~5)
\vspace{-12pt}
\end{table}

The results showed that participants systematically produced resumptive RCs even when they are disallowed in both the L1 and the target language. The L1-Greek and L1-Persian participants produced more resumptive RCs on average than the L1-Finnish and L1-Spanish participants, but importantly in each group there were at least some who produced resumptive RCs consistently (by the `≥ 80\% of trials' benchmark used in this dissertation) in one or more conditions. Some participants extended resumptive RCs into environments where they are not permitted in the L1. The overall production patterns for gapping and resumption aligned very closely with the NPAH. That is to say, gapping tended to be used in positions toward the top of the hierarchy and resumption tended to be used in positions toward the bottom of the hierarchy. 

While Hyltenstam for the most part describes the gapping vs. resumption opposition in terms of typological markedness, he makes clear that he thinks the real driver of the amped up rates of resumption in his data is not markedness \textit{per se}, but rather the processing asymmetries that underlie the NPAH:

\begin{quote}
I do not see typological markedness as an explanatory concept for structural regularities in second language acquisition data, where I would rather look for explanations at a processing level. \dots{} If it is the case, as seems reasonable, that the NP Accessibility Hierarchy has a psychological validity in the sense that relative clauses with functions low in the hierarchy are more difficult to process, then the use of pronominal copies can be seen as a means of reducing processing load where it is largest.\par\nobreak
\hfill (\cite{hyltenstam1990}, pp. 33--34)
\end{quote}
\vspace{-12pt}

\noindent In light of this insight, I attribute to Hyltenstam the idea that processing difficulty in real-time production of RCs can give rise to L2 resumption. 

A number of other researchers have examined L2 resumptive RCs in languages thought to lack grammatical RC resumption (mostly focusing on L2 English) since the early forays by \textcite{gass1979} and \textcite{hyltenstam1984,hyltenstam1990}, as shown in Table~\ref{tab:acq-study}. Studies in this line of research have consistently found that L2ers from all language backgrounds make at least some use of resumption and that rates of both (a)~resumptive RC production (\cite{gass1979,hyltenstam1984,hyltenstam1990,maghrabi1997,marefat.abdollahnejad2014,pavesi1986,tezel1998}) and (b)~resumptive RC acceptance (\cite{maghrabi1997,marefat.abdollahnejad2014,solaimani.etal2023,tezel1998}) are higher in syntactic environments which are believed to be more difficult for relativization. Furthermore, these studies have found that L2 performance is at least partly dependent on the L1 (\cite{alosaimi2021,gass1979,hyltenstam1984,hyltenstam1990,maghrabi1997,perpinan2020,solaimani.etal2023,tezel1998}) and that more advanced L2ers typically show a closer approximation to native-speaker performance (\cite{alosaimi2021,maghrabi1997,marefat.abdollahnejad2014,solaimani.etal2023}).

\begin{table}
\caption{Experimental studies on L2 resumptive RCs where target language lacks grammatical resumption}\label{tab:acq-study}
\raggedright\setstretch{1}
\begin{tblr}{colspec={lX}}\toprule
Study & Details \\\midrule
\cite{pavesi1986} & L1: Italian; L2: English; Method: oral elicitation; Environments: SU, DO, IO, OBL, GEN, OCOMP (short-distance) \\
\cite{maghrabi1997} & L1: Arabic; L2: English; Method: acceptability judgment, written elicitation, picture description; Environments: SU, DO, IO, OBL, GEN, OCOMP (short-distance) \\
\cite{tezel1998} & L1: Arabic, Chinese, Turkish; L2: English; Method: acceptability judgment, written elicitation; Environments: SU, DO, OBL (short-distance) \\
\cite{marefat.abdollahnejad2014} & L1: Persian; L2: English; Method: acceptability judgment, written translation; Environments: SU, DO, OBL (short-distance) \\
\cite{perpinan2020} & L1: Arabic, English; L2: Spanish; Method: acceptability judgment, written elicitation; Environments: OBL (short-distance vs. adjunct-island) \\
\cite{alosaimi2021} & L1: Arabic; L2: English; Method: acceptability judgment; Environments: SU, DO, OBL, OCOMP (short-distance vs. RC-island) \\
\cite{solaimani.etal2023} & L1: French, Persian; L2: English; Method: acceptability judgment; Environments: SU, DO, OBL (short-distance with \textit{who} vs. \textit{that}) \\
\bottomrule
\end{tblr}
\vspace{-12pt}
\end{table}

%========================================================%
\section{R. Hawkins and Chan (1997)}\label{acq-hawkins}
%========================================================%

The experiment reported by \textcite{hawkins.chan1997} merits inclusion in the present discussion, even though it was not specifically designed to examine resumptive RCs. First off, it is one of the relatively few L2 RC studies that does not focus exclusively on short-distance RCs. Second, it provides fairly strong evidence that resumptive RCs can be transferred from the L1 to the interlanguage at the beginning stages of L2 acquisition: Whereas \textcite{hyltenstam1984,hyltenstam1990} based his arguments regarding L1 transfer entirely on production data, R. Hawkins and Chan used a combination of judgment data and correction data.

Their main instrument was an AJT aimed at testing the implicit knowledge that L1-Cantonese L2ers of English (CLEs) and L1-French L2ers of English (FLEs) have of island constraints for English RCs. The design of this task is summarized in Table \ref{tblr:hawkins}.

\begin{longtblr}[caption={\textcite{hawkins.chan1997}: Acceptability judgment task}, label={tblr:hawkins}]{colspec={|X|}, rowhead=0, rows={font=\raggedright\setstretch{1}}}\hline
\textbf{Participants:} 292~participants divided among 7~groups: 147~L1-Cantonese L2ers of English aged 12--21 (47~elementary, 46~intermediate, 54~advanced), 113~L1-French L2ers of English aged 12--19 (33~elementary, 40~intermediate, 40~advanced), and 32~L1-English controls aged 16--19; English proficiency was assessed via the Oxford Placement Test (\cite{allan1992}) \vspace{5pt}\rule{0pt}{10pt}\\*
\textbf{Stimuli:} 101 stimulus sentences, of which 17 had RPs in single-clause RCs (SU:~2~tokens; DO:~3~tokens; IO:~3~tokens; OBL:~6~tokens; OCOMP:~3~tokens) and 8 had island-constraint violations (\textit{wh}-island:~4~tokens; other complex NP island:~4~tokens) \vspace{5pt}\\
\textbf{Procedure:} Participants rated sentences on a 4-point scale ranging from A (`definitely correct') to D (`definitely incorrect'); if they gave a low rating, they were asked to provide a correction \vspace{5pt}\\
\textbf{Relevant conditions (pp. 225--226, Appendices B and C):} \vspace{5pt}\\*
\begin{varwidth}{\linewidth}\begin{enumerate}[\bfseries (a),itemsep=5pt] 
\item \textbf{Resumption at SU:}\par\nobreak
*The waiter [that \boxing{he} always serves us] is called George.
\item \textbf{Resumption at DO:}\par\nobreak
*The patient [that I visited \boxing{him}] is very sick.
\item \textbf{Resumption at IO:}\par\nobreak
*The schoolboy [that I read a story to \boxing{him}] finished his homework.
\item \textbf{Resumption at OBL:}\par\nobreak
*The file [that she put the papers in \boxing{it}] has been stolen.
\item \textbf{Resumption at OCOMP:}\par\nobreak
*The sailor [that Bill is more experienced than \boxing{him}] owns a boat.
\item \textbf{\textit{Wh}-island ORC:}\par\nobreak
*This is the man [whom Mary told me [when she will visit \_\_]].
\item \textbf{Complex NP island ORC:}\par\nobreak
*This is the secretary [whom Peter heard the news [that the boss will marry \_\_]].
\end{enumerate}\end{varwidth}\vspace{5pt}\\*
\textit{Note}. Other conditions examined various grammatical and ungrammatical gapped RCs.
\\\hline
\end{longtblr}

\noindent Crucially, Cantonese is considered a grammatical resumption language whereas English and French are not. Table \ref{tab:hawkins} gives the distributions of gapping and resumption for these three languages in single-clause RCs as reported in \textcite{keenan.comrie1977}. While Keenan and Comrie do not present data on Cantonese, \textcite{hawkins.chan1997} claim that the facts in Cantonese are the same as in Mandarin (p.~195).

\begin{table}
\caption{Distribution of gapping (−), resumption (+), and unrelativizable domains (0) in single-clause RCs for the languages in \textcite{hawkins.chan1997}}\label{tab:hawkins}
\raggedright\setstretch{1}
\begin{tabularx}{\linewidth}{lZZZZZZ}\hline
Language & SU & DO & IO & OBL & GEN & OCOMP \\\hline
English & − & − & − & − & − & − \\
French & − & − & − & − & − & 0 \\
Cantonese & − & ± & + & + & + & + \\\hline
\end{tabularx}\par\nobreak\vspace{5pt}
\hfill (adapted from \cite{keenan.comrie1977}, p.~76, Table~1)
\vspace{-12pt}
\end{table}

The CLE results are of particular interest. To begin with, nearly half of the elementary CLEs (22/47~= 47\%) consistently accepted RP trials, marking them as `correct' at least 73\% of the time. For comparison, FLEs rejected\footnotemark{} RP trials at high rates across all proficiency levels (elementary: 81\%, intermediate: 96\%, advanced: 98\%). In the elementary CLE group, furthermore, the vast majority of consistent RP acceptors (18/22~= 82\%) `corrected' island trials by inserting an RP at least 75\% of the time. By contrast, CLEs who did not accept RP trials at high rates only rarely inserted RPs in island trials. R.~Hawkins and Chan also reported that rates of rejection for RP trials among CLEs increased dramatically across proficiency levels (elementary:~38\%, intermediate:~55\%, advanced:~90\%). These results were taken as evidence that CLEs initially consider resumption a licit option for relativization in their interlanguage English due to L1 transfer and that they are able develop sensitivity to the ill-formedness of resumptive RCs in English as proficiency increases.

\footnotetext{For the FLE group, \textcite{hawkins.chan1997} report the rate of RP rejection but not the rate of RP acceptance; it is not necessarily safe to assume that the acceptance rate was the inverse of the rejection rate.}

%========================================================%
\section{An underresearched processing hierarchy in L2 experimental research}\label{acq-slw}
%========================================================%

Most of the existing experimental work on L2 resumption has focused on the NPAH. There has been comparatively little L2 research on other sets of syntactic environments that are thought to differ in processing difficulty, such as the three-way contrast between short-distance, long-distance, and \textit{wh}-island RCs examined in this dissertation (see~Sections~\ref{gap-sl}--\ref{gap-lw}). I follow, e.g.,~\textcite{hammerly2022}, \textcite{hawkins2004}, \textcite{kluender.kutas1993}, and \textcite{ogrady2022a} in treating these environments as forming a hierarchy of difficulty for the processing of A-bar dependencies, as illustrated in (\ref{ex-slw-ogrady}).\footnotemark{}

\footnotetext{In (\ref{ex-slw-ogrady}) we see that \textcite{ogrady2022a} uses language specific to \textit{wh}-dependencies, but the same hierarchy of processing difficulty also applies to other types of A-bar dependencies, including RC dependencies.}

\pex 
Affordability metric for extending \textit{wh} dependencies:\label{ex-slw-ogrady}
\a no extension across a clause boundary
\a extension across a clause boundary if there is no other \textit{wh} dependency
\a extension across a clause boundary if there is another \textit{wh} dependency\par\nobreak\vspace{5pt}
\hfill (\cite{ogrady2022a}, p. 73)
\xe 

\noindent For ease of reference, I refer to this hierarchy as the `SLW hierarchy'\footnotemark{}: The `S' level corresponds to short-distance RCs, the `L' level corresponds to long-distance RCs, and the `W' level corresponds to \textit{wh}-island RCs. What sets L apart from S is that it requires the dependency to cross a clause boundary, and what sets W apart from L (in languages with overt \textit{wh}-movement) is that it requires the dependency to be extended into an embedded clause that has another filler waiting to be integrated (see~Sections~\ref{pro-routine}, \ref{gap-sl}, and \ref{gap-lw}). As shown in (\ref{ex-slw-scale}), English short-distance RCs are relatively easy to process, with the dependency resolving in the RC's main (i.e.,~highest or only) clause. English long-distance RCs up the ante by extending the dependency into a declarative complement clause that is the next level of embedding, and English \textit{wh}-island RCs further burden the processor by again extending the dependency into the next level of embedding, but this time instead into an interrogative complement clause with a fronted \textit{wh}-phrase; note that this \textit{wh}-phrase introduces a second gapped dependency that also needs to be resolved in that clause (\cite{kluender1998}, p.~265). \textit{Wh}-island RCs are typically regarded as ungrammatical in English, but as we saw in Section~\ref{gap-lw}, the processing asymmetry between L and W holds even for languages like Hebrew that allow relativization at all three levels on the hierarchy (\cite{keshev.meltzer-asscher2019}).

\footnotetext{I use `hierarchy' here in the sense of \textcite{hawkins2004}: Whenever there is a processing asymmetry A~>~B and another processing asymmetry B~>~C, then the sequence A~>~B~>~C can be said to form a hierarchy (p.~6).}

\pex\label{ex-slw-scale}
The SLW hierarchy of processing difficulty in gapped dependencies
\a Short-distance (`S'):\vspace{5pt}\newline
the man [that the officers arrested \_\_] >
\a Long-distance (`L'):\vspace{5pt}\newline
the man [that I think [the officers arrested \_\_]] >
\a \textit{Wh}-island (`W'):\vspace{5pt}\newline
\ljudge*the man [that I wonder [which officers arrested \_\_]]
\xe

%========================================================%
\section{The Performance-Competence Correspondence Assumption}\label{acq-assumption}
%========================================================%

An assumption, explicit or implicit, often made in the traditional L2 research literature, including by \textcite{hyltenstam1984,hyltenstam1990}, is that an L2er's production patterns generally reflect representations licensed by the interlanguage grammar. For instance, if an L2er produces RCs like the one in (\ref{ex-assumption}), researchers typically concluded that the L2er's interlanguage grammar allows resumption in genitive RCs. 

\ex\label{ex-assumption}
\ljudge*the man [that I met \boxing{his} mother]
\xe

This assumption, which I dub the `Performance-Competence Correspondence Assumption' (PCCA), is widespread in the L2 acquisition literature, particularly when it comes to L2 resumption in RCs. The quotations below provide examples of researchers drawing conclusions about the state of the interlanguage grammar based on production patterns in RC elicitation tasks. Boxes mark text to which I would like to draw the reader's attention.

\begin{quote}\begin{samepage}
The Spanish group would seem to be particularly interesting \dots~since in spite of the fact that Spanish, according to our description, does not use pronominal copies, such a device is employed to a considerable extent in their \ilb{Swedish interlanguage}.\par\nobreak\vspace{5pt}
\hfill(\cite{hyltenstam1990}, p. 30)\end{samepage}\end{quote}\vspace{-12pt}

\begin{quote}\begin{samepage}
The results obtained in this study support the general hypothesis that [L2 acquisition] progresses from unmarked to marked both in terms of the order in which \ilb{the various target language structures are mastered} and in terms of the different strategies employed in the process of acquisition.\par\nobreak\vspace{5pt}
\hfill(\cite{pavesi1986}, p. 49)        
\end{samepage}\end{quote}\vspace{-12pt}

\begin{quote}\begin{samepage}
The most interesting fact about \dots{} Hyltenstam's (\citeyear{hyltenstam1984}) study is that the L2 pattern of resumptive pronouns in relative clauses is found in neither the [native~language] nor the [target~language], but is represented in the grammars of a number of languages throughout the world. How could these learners have \ilb{concocted this [interlanguage]} \ilb{structure}, and \dots{} done so in a way that conforms to what occurs in other languages?
\par\nobreak\vspace{5pt}
\hfill(\cite{eckman2004}, p. 107)        
\end{samepage}\end{quote}\vspace{-12pt}

\begin{quote}\begin{samepage}
For the most part \dots~even when the participants performed in neither the target-like manner nor a native-like manner, \ilb{their [interlanguage] grammars} show greater tendency to provide resumption with types of relative clauses at the higher positions of the NPAH scale when they also provide resumption within the lower positions.
\par\nobreak\vspace{5pt}
\hfill(\cite{algady2013}, p. 184)        
\end{samepage}\end{quote}\vspace{-12pt}

\begin{quote}\begin{samepage}
The results of the translation of object-of-preposition RCs showed that most of the participants were either uncertain about using RPs or \ilb{considered them required}[,] which can be the result of direct L1 transfer. But at the advanced level, 93.3\% \ilb{considered RPs forbidden}.
\par\nobreak\vspace{5pt}
\hfill(\cite{marefat.abdollahnejad2014}, p. 36)        
\end{samepage}\end{quote}\vspace{-12pt}

In the wider linguistics community, the observation that production patterns do not always perfectly match the speaker's implicit knowledge goes back at least as far as the 1960s (\cite{chomsky1965}, p.~4; see~also \cite{kim2016}, p.~4ff; \cite{lust2006}, p.~124; \cite{schwartz1986}, p.~122ff; \cite{schwartz1993}, p.~150). Within the literature on resumption in RCs, a challenge to the PCCA comes from the experimental work on, e.g., L1-English adults reviewed in Chapter \ref{ch-resumption}. In acceptability judgment tasks, resumption trials tend to receive low ratings no matter the syntactic environment, indicating that they are considered ill-formed. In elicited production tasks, although participants generally do not produce resumptive RCs, when they do it tends to be at higher rates as the syntactic environment gets more complex, both in non-island contexts, where gapped RCs are well-formed, and in island contexts, where gapped RCs are, according to AJT data, ill-formed. Together these results suggest that resumptive RCs produced by L1-English adults are processing resumptives rather than grammatical resumptives. That is to say, they are ill-formed but still produced from time to time as a means of reducing processing load. 

The experiments included in the present dissertation test whether resumptive RCs produced by L2ers of English might also be processing resumptives, at least in some cases, contrary to the PCCA. If there is evidence that English resumptive RCs are strictly processing resumptives for some L2ers (i.e.,~they reject all conditions with resumptive RCs but still derive a processing benefit from resumption in RC production and/or RC comprehension), this would suggest that for those L2ers, discrepancies in performance between them and native speakers stem from differences in processing efficiency rather than differences in implicit knowledge. Sentence processing in general, after all, is known to require more effort in a nonnative language than in a native language (\cite{hopp2014,kilborn1992}). 

%========================================================%
\section{A preliminary L1/L2 acquisition pathway for single-clause RCs (to be revised)}\label{acq-preliminary}
%========================================================%

Let us now consider how L1/L2 learners might develop implicit knowledge about the distribution of well-formed gapped and resumptive RCs in English. I begin by piecing together an approach to the L1/L2 acquisition of single-clause RCs based on arguments presented in \textcite{hyltenstam1984,hyltenstam1990}, supplemented by statements made in Eckman's (\citeyear{eckman2004}) like-minded reanalysis of the same dataset. After fleshing out the details of this approach, I go on to point out some of the most glaring problems with it.

Hyltenstam and Eckman both leaned heavily into the terminology of typological markedness when presenting their arguments. As we have already established in this chapter, however, Hyltenstam saw processing rather than markedness as the root cause behind the emergence of resumptive RCs in interlanguages, at least when it cannot be straightforwardly traced to L1 transfer (see~also \cite{eckman2004}, p.~98ff). I therefore rephrase their arguments here so that they make appeal to processing considerations rather than to markedness. The key claims that Hyltenstam and Eckman make about the acquisition of single-clause RCs are listed in (\ref{ex-hyltenstam-pathway}).

\pex\label{ex-hyltenstam-pathway} Hyltenstam and Eckman's claims regarding the L1/L2 acquisition of RCs:\par\nobreak
\a For reasons of economy, when learners have a choice between something that is relatively easy on the processor and something that is relatively hard on the processor, they start by positing that only the easier option is permitted (\cite{eckman2004}, p.~98; \cite{hyltenstam1984}, p.~43)
\a On the assumption that resumption is easier on the processor than gapping is, the initial state for both L1 and L2 acquisition of single-clause RCs is to use only resumption across the entire NPAH (\cite{eckman2004}, p.~96; \cite{hyltenstam1984}, p.~45)
\a Over time, learners restructure the grammar, progressing from the initial state toward the target state, replacing resumption with gapping as necessary in response to positive evidence that gapping is available at the relevant position (\cite{eckman2004}, p.~104; \cite{hyltenstam1984}, p.~55)
\a The availability of gapping in one position on the hierarchy entails its availability in all higher positions (\cite{eckman2004}, p.~105; \cite{hyltenstam1984}, p.~45; see~\cite{keenan.comrie1977}, p.~68)
\xe

\noindent With these claims in mind, we can imagine the acquisition process unfolding in the manner displayed in Figure~\ref{fig:hyltenstam-pathway}. Recall that the position adopted here is that the NPAH represents a scale of processing difficulty with SU (easiest) at its top and OCOMP (hardest) at its bottom. Plus signs mark positions where resumption is permitted, minus signs mark positions where gapping is permitted, and square brackets mark relativization options that are backed up by positive evidence in the input. As we will see shortly, a crucial shortcoming of this approach is that it does not leave room for optionality between gapping and resumption in one or more positions on the hierarchy.

\begin{figure}
\caption{Preliminary L1/L2 acquisition pathway for single-clause RCs}
\label{fig:hyltenstam-pathway}
\setstretch{1}
\raggedright
\begin{mybox}
{\bfseries Step 0: Initial state}\par\vspace{2pt}
Assumed initial state has resumptive RCs across the entire scale
\par\vspace{6pt}
\tikz{
    \draw[draw=white, line width=1pt] (0,-5pt) -- (0,5pt);
    \draw[line width=1pt] (0.4,0) -- (12.4,0);
    \draw[line width=1pt] (0.4,-5pt) -- (0.4,5pt);
    \draw[line width=1pt] (2.4,-5pt) -- (2.4,5pt);
    \draw[line width=1pt] (4.4,-5pt) -- (4.4,5pt);
    \draw[line width=1pt] (6.4,-5pt) -- (6.4,5pt);
    \draw[line width=1pt] (8.4,-5pt) -- (8.4,5pt);
    \draw[line width=1pt] (10.4,-5pt) -- (10.4,5pt);
    \draw[line width=1pt] (12.4,-5pt) -- (12.4,5pt);
    \path(1.4,-.3)node(){\strut\ttfamily SU};
    \path(3.4,-.3)node(){\strut\ttfamily DO};
    \path(5.4,-.3)node(){\strut\ttfamily IO};
    \path(7.4,-.3)node(){\strut\ttfamily OBL};
    \path(9.4,-.3)node(){\strut\ttfamily GEN};
    \path(11.4,-.3)node(){\strut\ttfamily OCOMP};
    \path(1.4,.3)node(){\strut\large\ttfamily +};
    \path(3.4,.3)node(){\strut\large\ttfamily +};
    \path(5.4,.3)node(){\strut\large\ttfamily +};
    \path(7.4,.3)node(){\strut\large\ttfamily +};
    \path(9.4,.3)node(){\strut\large\ttfamily +};
    \path(11.4,.3)node(){\strut\large\ttfamily +};
}\par
\vspace{3pt}
{\bfseries Step 1: Positive evidence of gapped RCs at SU}\par\vspace{2pt}
Resumption replaces gapping at SU \par
\vspace{6pt}
\tikz{
    \draw[draw=white, line width=1pt] (0,-5pt) -- (0,5pt);
    \draw[line width=1pt] (0.4,0) -- (12.4,0);
    \draw[line width=1pt] (0.4,-5pt) -- (0.4,5pt);
    \draw[line width=1pt] (2.4,-5pt) -- (2.4,5pt);
    \draw[line width=1pt] (4.4,-5pt) -- (4.4,5pt);
    \draw[line width=1pt] (6.4,-5pt) -- (6.4,5pt);
    \draw[line width=1pt] (8.4,-5pt) -- (8.4,5pt);
    \draw[line width=1pt] (10.4,-5pt) -- (10.4,5pt);
    \draw[line width=1pt] (12.4,-5pt) -- (12.4,5pt);
    \path(1.4,-.3)node(){\strut\ttfamily SU};
    \path(3.4,-.3)node(){\strut\ttfamily DO};
    \path(5.4,-.3)node(){\strut\ttfamily IO};
    \path(7.4,-.3)node(){\strut\ttfamily OBL};
    \path(9.4,-.3)node(){\strut\ttfamily GEN};
    \path(11.4,-.3)node(){\strut\ttfamily OCOMP};
    \path(1.4,.3)node(){\strut\large\ttfamily [-]};
    \path(3.4,.3)node(){\strut\large\ttfamily +};
    \path(5.4,.3)node(){\strut\large\ttfamily +};
    \path(7.4,.3)node(){\strut\large\ttfamily +};
    \path(9.4,.3)node(){\strut\large\ttfamily +};
    \path(11.4,.3)node(){\strut\large\ttfamily +};
}\par
\vspace{3pt}
{\bfseries Step 2: Positive evidence of gapped RCs at IO}\par\vspace{2pt}
Resumption replaces gapping at IO and all higher levels on the scale \par
\vspace{6pt}
\tikz{
    \draw[draw=white, line width=1pt] (0,-5pt) -- (0,5pt);
    \draw[line width=1pt] (0.4,0) -- (12.4,0);
    \draw[line width=1pt] (0.4,-5pt) -- (0.4,5pt);
    \draw[line width=1pt] (2.4,-5pt) -- (2.4,5pt);
    \draw[line width=1pt] (4.4,-5pt) -- (4.4,5pt);
    \draw[line width=1pt] (6.4,-5pt) -- (6.4,5pt);
    \draw[line width=1pt] (8.4,-5pt) -- (8.4,5pt);
    \draw[line width=1pt] (10.4,-5pt) -- (10.4,5pt);
    \draw[line width=1pt] (12.4,-5pt) -- (12.4,5pt);
    \path(1.4,-.3)node(){\strut\ttfamily SU};
    \path(3.4,-.3)node(){\strut\ttfamily DO};
    \path(5.4,-.3)node(){\strut\ttfamily IO};
    \path(7.4,-.3)node(){\strut\ttfamily OBL};
    \path(9.4,-.3)node(){\strut\ttfamily GEN};
    \path(11.4,-.3)node(){\strut\ttfamily OCOMP};
    \path(1.4,.3)node(){\strut\large\ttfamily [−]};
    \path(3.4,.3)node(){\strut\large\ttfamily −};
    \path(5.4,.3)node(){\strut\large\ttfamily [−]};
    \path(7.4,.3)node(){\strut\large\ttfamily +};
    \path(9.4,.3)node(){\strut\large\ttfamily +};
    \path(11.4,.3)node(){\strut\large\ttfamily +};
}\par
\vspace{3pt}
{\bfseries Step 3: Positive evidence of gapped RCs at OCOMP}
\par\vspace{2pt}
Resumption replaces gapping at OCOMP and all higher levels on the scale
\par\vspace{6pt}
\tikz{
    \draw[draw=white, line width=1pt] (0,-5pt) -- (0,5pt);
    \draw[line width=1pt] (0.4,0) -- (12.4,0);
    \draw[line width=1pt] (0.4,-5pt) -- (0.4,5pt);
    \draw[line width=1pt] (2.4,-5pt) -- (2.4,5pt);
    \draw[line width=1pt] (4.4,-5pt) -- (4.4,5pt);
    \draw[line width=1pt] (6.4,-5pt) -- (6.4,5pt);
    \draw[line width=1pt] (8.4,-5pt) -- (8.4,5pt);
    \draw[line width=1pt] (10.4,-5pt) -- (10.4,5pt);
    \draw[line width=1pt] (12.4,-5pt) -- (12.4,5pt);
    \path(1.4,-.3)node(){\strut\ttfamily SU};
    \path(3.4,-.3)node(){\strut\ttfamily DO};
    \path(5.4,-.3)node(){\strut\ttfamily IO};
    \path(7.4,-.3)node(){\strut\ttfamily OBL};
    \path(9.4,-.3)node(){\strut\ttfamily GEN};
    \path(11.4,-.3)node(){\strut\ttfamily OCOMP};
    \path(1.4,.3)node(){\strut\large\ttfamily [−]};
    \path(3.4,.3)node(){\strut\large\ttfamily −};
    \path(5.4,.3)node(){\strut\large\ttfamily [−]};
    \path(7.4,.3)node(){\strut\large\ttfamily −};
    \path(9.4,.3)node(){\strut\large\ttfamily −};
    \path(11.4,.3)node(){\strut\large\ttfamily [−]};
}
\vspace{-6pt}
\end{mybox}
\vspace{-6pt}
\textit{Note}. − = gapping; + = resumption; [ ] = direct evidence.\par 
\setstretch{1.5}\vspace{-12pt}
\end{figure}

Step~0 shows the L1/L2 initial state for the acquisition of single-clause RCs. It has mandatory resumption across the entire scale, following Hyltenstam and Eckman's assumption that resumption is the default option for relativization. In each of the steps that follow, the learner responds to positive evidence of gapped RCs at a given position by replacing resumption with gapping there as well as in all higher positions. I do not comment here on the quantity or quality of input that is required to initiate the restructuring process (for discussion of that topic, see~\cite{omaki.lidz2015}). Also, note that I am not making any specific claims about the order in which learners are exposed to RC input of different types; instead, I am merely offering an example of how the relevant input could be sequenced so that I can demonstrate how the model operates. In fact, I expect it is often the case that different learners are exposed to key pieces of positive evidence in different orders.

Step~1 shows how the learner responds to positive evidence of gapped RCs at SU (e.g.,~\textit{the man that hired me}). If the input is sufficient to trigger restructuring, the learner replaces mandatory resumption with mandatory gapping in that position. Notice that this assumes that gapping and resumption are in complementary distribution, which is the case for languages like English and Korean but not for languages like Hebrew and Mandarin.

Step~2 shows that the learner responds to positive evidence of gapping at IO (e.g.,~\textit{the man I gave the money to}) by extending gapping to both IO and DO. Note that independent evidence of gapping at DO is unnecessary because of the assumption made by both Hyltenstam and Eckman, stemming from the cross-linguistic work of \textcite{keenan.comrie1977}, that if gapping is permitted at one position on the NPAH, it is also permitted at all higher positions. Finally, Step~3 shows that the learner responds to positive evidence of gapping at OCOMP (e.g.,~\textit{the man I am taller than}) by extending gapping down to the bottom of the scale, resulting in the target grammar state with mandatory gapping across the entire NPAH.

As I see it, the proposal outlined above has two main strengths. The first is its simplicity. All that is required is a hierarchy of processing difficulty and positive evidence of the kind that we can reasonably expect learners to encounter in the input. The other strength of this proposal is each grammar state that learners pass through on their way to the target state aligns well with experimental data on the acquisition of RCs, both in L2 and in L2. Starting with the latter, the proposal is consistent with, e.g., Hyltenstam's finding that with respect to the NPAH, L2ers tended to produce gapped RCs at higher positions and resumptive RCs at lower positions in their acquisition of Swedish, a language without resumptive RCs. The proposal is also consistent with findings in the L1 acquisition literature: For instance, in a series of RC elicitation experiments, \textcite{kim2013} found that L1-English children aged 5--6 resorted to resumption at DO, IO, and OBL but not at SU; moreover, the proportion of resumptive RC responses increased down the NPAH from SU (0\%) to IO (11.4\%) and---with a different group of participants---from DO (5.7\%) to OBL (29.5\%).

There are also four important problems with this proposal, all of which stem from treating resumption as the default option for relativization. As we saw in Section~\ref{res-gapping}, few if any languages permit resumption at SU in short-distance RCs, and there is strong psycholinguistic evidence that the processor often prefers gapping over resumption in easy-to-process RCs. Second, positing an initial state with resumption at all levels on the hierarchy leaves no room for states in which neither gapping nor resumption is permitted at certain positions. That is to say, there is no mechanism for removing resumption from the hierarchy aside from replacing it with gapping. This is a critical limitation because there are many languages where relativization is not permitted at all for one or more grammatical functions on the NPAH. A few examples from \textcite{keenan.comrie1977} include Finnish, German, Korean, Malagasy, Spanish, and Tagalog (pp.~76--79). Third, if positive evidence of gapping (somehow) leads to the replacement of mandatory resumption with mandatory gapping at the relevant level on the scale, then it is unclear how optionality between gapped and resumptive RCs of the kind that Keenan and Comrie observed in languages like Mandarin and Persian can ever arise. Finally, although there are experimental studies showing that L1 children produce nontarget resumptive RCs at elevated rates compared to L1 adults (e.g.,~\cite{kim2013,labelle1990,perez-leroux1995,varlokosta.armon-lotem1998}), there is no evidence that L1 children in any language start off exclusively with resumptive RCs.

Before moving on, it bears mentioning that Hyltenstam equivocated somewhat on the matter of the L2 initial state. Towards the beginning of his seminal \citeyear{hyltenstam1984} paper, he clearly stated that resumption should be the default at every level on the hierarchy (p.~45). After observing that his Greek and Persian participants produced more resumptive RCs than his Finnish and Spanish ones did, however, he conceded that L2 initial states may vary based on L1 background, saying:

\begin{quote}\begin{samepage}
It is very likely that learners start at different points in the hierarchy depending on the nature of their mother tongue. Finnish learners [of an additional language], for example, probably never use pronominal copies in higher positions.
\par\nobreak\vspace{5pt}
\hfill(\cite{hyltenstam1984}, p. 55)\end{samepage}\end{quote}\vspace{-12pt}

What Hyltenstam seems to have in mind here is an L2 initial state that is somewhere between his assumed default L1 initial state (which has mandatory resumption across the NPAH) and the L2er's current grammar state in the L1 (which presumably has gapping in at least some positions on the NPAH). Depending on how far along the NPAH gapping is in the L1 grammar, the L2er may regress all the way to the assumed default L1 initial state or only part of the way there. However, Hyltenstam makes no attempt to explain how far exactly L2ers would regress or what prevents them from going further. In the next two sections, I propose a revised approach to the development of implicit knowledge of the distribution of gapping vs. resumption in the L1 and L2 acquisition of RCs, taking care to avoid the pitfalls of this preliminary proposal.

%========================================================%
\section{Principles underpinning processing-guided approaches to language acquisition}\label{acq-principles}
%========================================================%

The years intervening between \textcite{hyltenstam1984,hyltenstam1990} and the present day have seen a growing interest in processing-guided approaches to language acquisition, both among generativists (e.g.,~\cite{chomsky2017}) and among researchers working in other paradigms (e.g.,~\cite{ogrady2005}). For present purposes, I define a `processing-guided' approach as one in which hierarchies of processing difficulty work in concert with positive evidence in the input and (learning) principles relevant to grammar building to shepherd learners toward the target grammar. Such approaches are appealing for a number of reasons. First, they draw on processing considerations, such as the subject-object asymmetry in the processing of RCs, for which there is independent psycholinguistic evidence (\cite{hawkins2004}, p.~9; \cite{ogrady2022a}, p.~60; \cite{yang.etal2017}, p.~116). Second, they reduce the explanatory burden placed on the grammar (\cite{chomsky2005}, p.~9; \cite{chomsky2017}, p.~296). Last but not least, they recast the learnability problem such that the positive input appears less impoverished than it did previously (\cite{hawkins2007}, p.~339; \cite{ogrady2005}, p.~166). 

This section introduces the reader to five principles (`Gatekeeping,' `Conservatism,' `Affordability,' `Continuity,' and `Preemption') that I see as necessary components of any approach to language acquisition in which the development of implicit grammatical knowledge is guided by processing considerations. The principles need not all be regarded as processing principles in a strict sense; for my part, I consider them to be guidelines for grammar construction that work in conjunction with hierarchies of processing difficulty to help learners reach the target grammar state. The reader will notice a partial overlap between these principles and the claims from both Hyltenstam and Eckman that were considered in the previous section. Each definition is accompanied by a few short quotations illustrating how the principle has been described by other authors. These principles form the basis of the revised L1 and L2 acquisition pathways for RCs presented in subsequent sections of the chapter.

\subsection{Gatekeeping}\label{acq-gatekeeping}

The first principle states that processing factors (in both production and comprehension) create a spectrum of difficulty punctuated by natural break points (\cite{hawkins2004}, p.~192ff; \cite{ogrady2005}, p.~203ff; \cite{ogrady2011}, p.~34; \cite{ogrady2022a}, p.~133). In a sense, these break points act like gatekeepers, playing an instrumental role in determining the distribution of patterns such as gapped and resumptive RCs within a language; I therefore call this principle `Gatekeeping.' Only processing routines that involve an added operation or other complication that results in a substantial uptick in cognitive load constitute a break point; being forced to process a low-frequency lexical item therefore does not amount to a break point, but needing to reactivate a filler in a complement clause does (\cite{ogrady2022a}, p.~73). On the assumption that the NPAH is at its core a scale of processing difficulty (\cite{hawkins1999}, p.~252ff; \cite{keenan.comrie1977}, p.~88ff), it could conceivably contain up to five such break points: one between SU and DO, another between DO and IO, etc. I am not certain that I personally endorse the idea that every single level on the NPAH constitutes a break point, but let us suppose that it does for the remainder of the chapter, since I find it useful for purposes of exposition to use a scale with a relatively large number of levels.\footnotemark{}

\footnotetext{Examples of studies that have proposed collapsing some of the levels on the NPAH (e.g.,~IO and OBL) include \textcite{hawkins1999}, \textcite{herrmann2005}, and \textcite{hogbin.song2007}.}

\begin{quote}
\begin{samepage}
\dots{} processing is easier the more reduced, syntactically and semantically, a filler gap domain is. \dots{} Reduction motivates hierarchies of gap-containing environments, with grammars cutting off and with progressively fewer exemplifying languages down each hierarchy. 
\par\nobreak\vspace{5pt}
\hfill (\cite{hawkins1999}, pp. 247--248)
\end{samepage}\\\vspace{12pt}

\begin{samepage}
\dots{} processing factors create a continuum of difficulty along which there are certain natural `break points.'
\par\nobreak\vspace{5pt}
\hfill (\cite{ogrady2011}, p. 34)
\end{samepage}
\end{quote}\vspace{-12pt}

\subsection{Conservatism}\label{acq-conservatism}

The second principle is a general principle of learning which, as concerns language, states that learners start by positing the smallest grammar/language possible and only expand upon it if forced to do so by positive evidence in the input (\cite{berwick1985}, p.~235ff; \cite{clark1992}, p.~101; \cite{fodor.sakas2005}, p.~517; \cite{pinker1994}, p.~282). In the context of RC acquisition, I~interpret this to mean that learners are conservative in the sense that they do not generalize a relativization option (gapping or resumption) from an environment where they know it to be permissible to an environment that places substantially more strain on the processor (one that lies across a natural break point in difficulty/complexity, i.e., to the right on the hierarchy) unless there is positive evidence for that option, too  (\cite{hawkins2004}, p.~273; \cite{keenan.comrie1977}, p.~88; \cite{ogrady2005}, p.~184; \cite{ogrady2022a}, p.~129ff). In the L1 acquisition of RCs, this principle prevents children from generalizing from gapping at SU (the easiest position) to gapping at DO (a harder position) without positive evidence that gapping is also possible there. I follow \textcite{ogrady2005} in calling this principle `Conservatism' (p.~184).

\begin{quote}
\begin{samepage}
\dots{} a well-designed child, when faced with several choices in how far to generalize, should, in general, be conservative: start with the smallest hypothesis about the language that is consistent with what parents say, then expand it outward as the evidence requires.
\par\nobreak\vspace{5pt}
\hfill (\cite{pinker1994}, p. 282)
\end{samepage}\\\vspace{12pt}

\begin{samepage}
\dots{} the simplest environments on each hierarchy are postulated first, on account of their processing ease, and any extension beyond them to lower hierarchy positions will be forced on the learner by positive evidence from the relevant language.
\par\nobreak\vspace{5pt}
\hfill (\cite{hawkins2004}, p. 273)    
\end{samepage}\\\vspace{12pt}

\begin{samepage}
In a very real sense, children do indeed know exactly what to do from the outset: as little as possible. \dots{} Adoption of a more costly option \dots{} comes later and only in response to patterns of usage observed in the speech of others.
\par\nobreak\vspace{5pt}
\hfill (\cite{ogrady2022a}, pp. 129--131)
\end{samepage}
\end{quote}\vspace{-12pt}

\subsection{Affordability}\label{acq-affordability}

The third principle is a principle of processing economy which states that if a language allows a difficult operation, it will also allow easier operations of the same type (\cite{hawkins2004}, p.~100; \cite{huang2000}, p.~93; \cite{keenan.comrie1977}, p.~88; \cite{ogrady2022a}, p.~130). For example, provided with positive evidence that relativization is possible at OBL, learners should assume that relativization is also possible at SU through IO. Furthermore, if one option for relativization (e.g.,~resumption) is available at OBL, then it should also be available at higher levels on the scale, up to where another option (e.g.,~gapping) is available. I follow \textcite{ogrady2022a} in referring to this principle as `Affordability' (p.~43). 

\begin{quote}
\begin{samepage}
\dots{} if speakers can in general encode a meaning of a certain difficulty in some way, then they should certainly be able to encode meanings of lesser difficulty (possibly in a different way).
\par\nobreak\vspace{5pt}
\hfill (\cite{keenan.comrie1977}, p. 88)
\end{samepage}\\\vspace{12pt}

\begin{samepage}
If a language permits a more costly operation, it should also permit a less costly operation of the same type.\par\nobreak\vspace{5pt}
\hfill (\cite{ogrady2022a}, p. 43)
\end{samepage}
\end{quote}\vspace{-12pt}

\subsection{Continuity}\label{acq-continuity}

The fourth principle is derived from Keenan and Comrie's (\citeyear{keenan.comrie1977}) observation that languages appear to require that each option for relativization operate on an uninterrupted stretch of the NPAH. Specifically, the principle states that if a way of encoding meaning is permitted at two distinct levels on the relevant hierarchy, then it will also be permitted at all intermediate levels. Learners who allow resumptive RCs at both DO and OBL should therefore also allow resumptive RCs at IO. I~refer to this principle as `Continuity.'

\begin{quote}
Any RC-forming strategy must apply to a continuous segment of the [NPAH].
\par\nobreak\vspace{5pt}
\hfill (\cite{keenan.comrie1977}, p. 67)
\end{quote}\vspace{-12pt}

% \begin{samepage}
% \dots{} if a given strategy is used to encode a fairly easy meaning and that strategy is ``strong'' enough to encode a rather difficult meaning, then it is surely strong enough to encode the meanings of intermediate difficulty.
% \par\nobreak\vspace{5pt}
% \hfill (\cite{keenan.comrie1977}, p. 77)
% \end{samepage}
% \end{quote}\vspace{-12pt}

\subsection{Preemption}\label{acq-preemption}

The fifth and final principle is a general principle of learning which, as concerns language, states that if one option for encoding a meaning is supported by positive evidence in the input, learners will initially assume that other options for encoding the same meaning are impermissible (\cite{berwick1985}; \cite{pinker1984}, p.~113; \cite{rutherford1986}, p.~443; \cite{trahey1992}, p.~31). The possibility that each option is allowed will not be entertained without it being forced. The way I see it, there are two things that can force an option to be present at a given position: (a)~direct evidence in the input that the option is available there or (b)~direct evidence in the input that the option is available both above and below the relevant position, such that Continuity demands that it be permitted there as well. For instance, given positive evidence of gapped RCs, learners should assume that the language allows only gapped RCs in the relevant environment unless they encounter positive evidence which---either directly or via Continuity---forces resumptive RCs to also be allowed there. I follow \textcite{rutherford1989} in calling this principle `Preemption' (p.~443). Another name that it goes by is the Uniqueness Principle (\cite{pinker1984}, p.~113).

\begin{quote}
\begin{samepage}
The principle is that the need for negative evidence in language acquisition can be eliminated if the child knows that when he or she is faced with a set of alternative structures fulfilling the same function, only one of the structures is correct unless there is direct evidence that more than one is necessary.
\par\nobreak\vspace{5pt}
\hfill (\cite{pinker1984}, p. 113)
\end{samepage}\\\vspace{12pt}

\begin{samepage}
[O]f two or more formal entries for encoding a semantic feature or representing a semantic relationship, one and only one is required. \dots{} [O]ne of two possible candidates for selection is preempted because there is room for only one and a selection \textit{must} [emphasis~original] be made.
\par\nobreak\vspace{5pt}
\hfill (\cite{rutherford1989}, p. 443)
\end{samepage}\\\vspace{12pt}

\begin{samepage}
In general, preemption ensures a one-to-one mapping such that a single meaning is mapped onto one morphological feature or syntactic relationship. When two different forms encoding the same meaning are in competition, only one can be retained. \dots{} It is assumed within the framework discussed here that the form attested in the input is the one that will be retained.
\par\nobreak\vspace{5pt}
\hfill (\cite{trahey1992}, p. 31)
\end{samepage}
\end{quote}\vspace{-12pt}

These principles should be compatible with any approach to linguistic analysis wherein processing pressures and positive evidence in the input both have roles to play in guiding language acquisition. Something I have not dealt with here is how learners discriminate between reliable input (well-formed expressions) and unreliable input (performance errors and other ill-formed expressions). For detailed discussion on that topic, see Yang's (\citeyear{yang2000}) dissertation.

%========================================================%
\section{A revised L1 acquisition pathway for single-clause RCs}\label{acq-revised}
%========================================================%

Figure \ref{fig:revised-pathway} shows a revised L1 acquisition pathway for single-clause RCs which draws on the five principles introduced in the previous section. The target language is a hypothetical language with mandatory gapping at SU, optionality between gapping and resumption at DO through IO, mandatory resumption at OBL through GEN, and no relativization at OCOMP. The reason that I elected to use a hypothetical language here rather than an actual language, such as English, is to showcase how the principles interact in a relatively complex test case. 

\begin{figure}
\caption{Revised L1 acquisition pathway for single-clause RCs (in a hypothetical language---not English)}
\label{fig:revised-pathway}
\setstretch{1}
\raggedright
\begin{mybox}
{\bfseries Step 0: Initial state}
\par\vspace{2pt}
Assumed initial state has no relativization anywhere on the scale
\par\vspace{6pt}
\tikz{
    \draw[draw=white, line width=1pt] (0,-5pt) -- (0,5pt);
    \draw[line width=1pt] (0.4,0) -- (12.4,0);
    \draw[line width=1pt] (0.4,-5pt) -- (0.4,5pt);
    \draw[line width=1pt] (2.4,-5pt) -- (2.4,5pt);
    \draw[line width=1pt] (4.4,-5pt) -- (4.4,5pt);
    \draw[line width=1pt] (6.4,-5pt) -- (6.4,5pt);
    \draw[line width=1pt] (8.4,-5pt) -- (8.4,5pt);
    \draw[line width=1pt] (10.4,-5pt) -- (10.4,5pt);
    \draw[line width=1pt] (12.4,-5pt) -- (12.4,5pt);
    \path(1.4,-.3)node(){\strut\ttfamily SU};
    \path(3.4,-.3)node(){\strut\ttfamily DO};
    \path(5.4,-.3)node(){\strut\ttfamily IO};
    \path(7.4,-.3)node(){\strut\ttfamily OBL};
    \path(9.4,-.3)node(){\strut\ttfamily GEN};
    \path(11.4,-.3)node(){\strut\ttfamily OCOMP};
    \path(1.4,.3)node(){\strut\ttfamily 0};
    \path(3.4,.3)node(){\strut\ttfamily 0};
    \path(5.4,.3)node(){\strut\ttfamily 0};
    \path(7.4,.3)node(){\strut\ttfamily 0};
    \path(9.4,.3)node(){\strut\ttfamily 0};
    \path(11.4,.3)node(){\strut\ttfamily 0};
}\par
\vspace{3pt}
{\bfseries Step 1: Positive evidence of gapped RCs at SU}
\par\vspace{2pt}
Gapping added at SU
\par\vspace{6pt}
\tikz{
    \draw[draw=white, line width=1pt] (0,-5pt) -- (0,5pt);
    \draw[line width=1pt] (0.4,0) -- (12.4,0);
    \draw[line width=1pt] (0.4,-5pt) -- (0.4,5pt);
    \draw[line width=1pt] (2.4,-5pt) -- (2.4,5pt);
    \draw[line width=1pt] (4.4,-5pt) -- (4.4,5pt);
    \draw[line width=1pt] (6.4,-5pt) -- (6.4,5pt);
    \draw[line width=1pt] (8.4,-5pt) -- (8.4,5pt);
    \draw[line width=1pt] (10.4,-5pt) -- (10.4,5pt);
    \draw[line width=1pt] (12.4,-5pt) -- (12.4,5pt);
    \path(1.4,-.3)node(){\strut\ttfamily SU};
    \path(3.4,-.3)node(){\strut\ttfamily DO};
    \path(5.4,-.3)node(){\strut\ttfamily IO};
    \path(7.4,-.3)node(){\strut\ttfamily OBL};
    \path(9.4,-.3)node(){\strut\ttfamily GEN};
    \path(11.4,-.3)node(){\strut\ttfamily OCOMP};
    \path(1.4,.3)node(){\strut\large\ttfamily [−]};
    \path(3.4,.3)node(){\strut\ttfamily 0};
    \path(5.4,.3)node(){\strut\ttfamily 0};
    \path(7.4,.3)node(){\strut\ttfamily 0};
    \path(9.4,.3)node(){\strut\ttfamily 0};
    \path(11.4,.3)node(){\strut\ttfamily 0};
}\par
\vspace{3pt}
{\bfseries Step 2: Positive evidence of gapped RCs at IO}
\par\vspace{2pt}
Gapping added at IO; Affordability demands that gapping be extended upward from IO
\par\vspace{6pt}
\tikz{
    \draw[draw=white, line width=1pt] (0,-5pt) -- (0,5pt);
    \draw[line width=1pt] (0.4,0) -- (12.4,0);
    \draw[line width=1pt] (0.4,-5pt) -- (0.4,5pt);
    \draw[line width=1pt] (2.4,-5pt) -- (2.4,5pt);
    \draw[line width=1pt] (4.4,-5pt) -- (4.4,5pt);
    \draw[line width=1pt] (6.4,-5pt) -- (6.4,5pt);
    \draw[line width=1pt] (8.4,-5pt) -- (8.4,5pt);
    \draw[line width=1pt] (10.4,-5pt) -- (10.4,5pt);
    \draw[line width=1pt] (12.4,-5pt) -- (12.4,5pt);
    \path(1.4,-.3)node(){\strut\ttfamily SU};
    \path(3.4,-.3)node(){\strut\ttfamily DO};
    \path(5.4,-.3)node(){\strut\ttfamily IO};
    \path(7.4,-.3)node(){\strut\ttfamily OBL};
    \path(9.4,-.3)node(){\strut\ttfamily GEN};
    \path(11.4,-.3)node(){\strut\ttfamily OCOMP};
    \path(1.4,.3)node(){\strut\large\ttfamily [−]};
    \path(3.4,.3)node(){\strut\large\ttfamily −};
    \path(5.4,.3)node(){\strut\large\ttfamily [−]};
    \path(7.4,.3)node(){\strut\ttfamily 0};
    \path(9.4,.3)node(){\strut\ttfamily 0};
    \path(11.4,.3)node(){\strut\ttfamily 0};
}\par
\vspace{3pt}
{\bfseries Step 3: Positive evidence of resumptive RCs at GEN} 
\par\vspace{2pt}
Resumption added at GEN; Affordability demands that resumption be extended upward from GEN; Preemption blocks upward extension of resumption at IO
\par\vspace{6pt}
\tikz{
    \draw[draw=white, line width=1pt] (0,-5pt) -- (0,5pt);
    \draw[line width=1pt] (0.4,0) -- (12.4,0);
    \draw[line width=1pt] (0.4,-5pt) -- (0.4,5pt);
    \draw[line width=1pt] (2.4,-5pt) -- (2.4,5pt);
    \draw[line width=1pt] (4.4,-5pt) -- (4.4,5pt);
    \draw[line width=1pt] (6.4,-5pt) -- (6.4,5pt);
    \draw[line width=1pt] (8.4,-5pt) -- (8.4,5pt);
    \draw[line width=1pt] (10.4,-5pt) -- (10.4,5pt);
    \draw[line width=1pt] (12.4,-5pt) -- (12.4,5pt);
    \path(1.4,-.3)node(){\strut\ttfamily SU};
    \path(3.4,-.3)node(){\strut\ttfamily DO};
    \path(5.4,-.3)node(){\strut\ttfamily IO};
    \path(7.4,-.3)node(){\strut\ttfamily OBL};
    \path(9.4,-.3)node(){\strut\ttfamily GEN};
    \path(11.4,-.3)node(){\strut\ttfamily OCOMP};
    \path(1.4,.3)node(){\strut\large\ttfamily [−]};
    \path(3.4,.3)node(){\strut\large\ttfamily −};
    \path(5.4,.3)node(){\strut\large\ttfamily [−]};
    \path(7.4,.3)node(){\strut\large\ttfamily +};
    \path(9.4,.3)node(){\strut\large\ttfamily [+]};
    \path(11.4,.3)node(){\strut\ttfamily 0};
}\par
\vspace{3pt}
{\bfseries Step 4: Positive evidence of resumptive RCs at DO} 
\par\vspace{2pt}
Resumption added at DO; Continuity demands (a)~that gapping be available at SU through IO and (b)~that resumption be available at DO though GEN
\par\vspace{6pt}
\tikz{
    \draw[draw=white, line width=1pt] (0,-5pt) -- (0,5pt);
    \draw[line width=1pt] (0.4,0) -- (12.4,0);
    \draw[line width=1pt] (0.4,-5pt) -- (0.4,5pt);
    \draw[line width=1pt] (2.4,-5pt) -- (2.4,5pt);
    \draw[line width=1pt] (4.4,-5pt) -- (4.4,5pt);
    \draw[line width=1pt] (6.4,-5pt) -- (6.4,5pt);
    \draw[line width=1pt] (8.4,-5pt) -- (8.4,5pt);
    \draw[line width=1pt] (10.4,-5pt) -- (10.4,5pt);
    \draw[line width=1pt] (12.4,-5pt) -- (12.4,5pt);
    \path(1.4,-.3)node(){\strut\ttfamily SU};
    \path(3.4,-.3)node(){\strut\ttfamily DO};
    \path(5.4,-.3)node(){\strut\ttfamily IO};
    \path(7.4,-.3)node(){\strut\ttfamily OBL};
    \path(9.4,-.3)node(){\strut\ttfamily GEN};
    \path(11.4,-.3)node(){\strut\ttfamily OCOMP};
    \path(1.4,.3)node(){\strut\large\ttfamily [−]};
    \path(3.4,.3)node(){\strut\large\ttfamily -/[+]};
    \path(5.4,.3)node(){\strut\large\ttfamily [-]/+};
    \path(7.4,.3)node(){\strut\large\ttfamily +};
    \path(9.4,.3)node(){\strut\large\ttfamily [+]};
    \path(11.4,.3)node(){\strut\ttfamily 0};
}
\vspace{-6pt}
\end{mybox}
\vspace{-6pt}
\textit{Note}. 0 = unrelativizable; − = gapping; + = resumption; [ ] = direct evidence.\par
\setstretch{1.5}\vspace{-12pt}
\end{figure}

In observance of Conservatism, the initial state at Step~0 has no relativization anywhere on the scale. Doing things this way allows for the possibility of final states where relativization is excluded altogether from one or more positions. Like in Figure~\ref{fig:hyltenstam-pathway}, the remaining steps show how learners respond in the model if they are next exposed to evidence of gapped RCs at SU, then at IO, and finally at OCOMP. As a reminder, I am not making any specific claims about the order in which learners are exposed to the relevant input. Instead, my intent is simply to illustrate how acquisition might proceed given one possible ordering.

At Step~1, the learner responds to positive evidence of gapped RCs at SU by restructuring the grammar to allow gapping at that position. Conservatism prevents the learner from generalizing gapping to any other positions, since they are all lower on the NPAH. 

At Step~2, the learner responds to positive evidence of gapped RCs at IO adding gapping there. Affordability demands that gapping then be extended upward from IO.  

At Step~3, the learner responds to positive evidence of resumptive RPs at GEN by adding resumption there. Affordability demands that resumption then be extended upward from GEN but only as far as IO because Preemption blocks resumption from extending up past OBL. Note that it is only possible to determine whether an option is forced to be present (by either positive evidence or Continuity, as the case may be) if learners subconsciously track and retain information regarding the environments for which they have encountered direct evidence of that option.

Finally, at Step~4, the learner responds to positive evidence of resumptive RCs at DO by adding resumption there. Because there has been positive evidence of gapped RCs at both SU and IO, Continuity demands that gapping also be retained at the intervening DO position. By the same token, because there has been positive evidence of resumptive RCs at both DO and GEN, Continuity demands that resumption also be available at the intervening IO and OBL positions. The consequence is optionality between gapping and resumption at DO and IO. Note that Preemption as defined in Section~\ref{acq-preemption} allows for optionality, but only in cases where each way of encoding the meaning is forced at the relevant position. The result is the target grammar state with mandatory gapping at SU, optionality between gapping and resumption at DO through IO, mandatory resumption at OBL through GEN, and no relativization at OCOMP. 

What happens, though, if the steps in Figure~\ref{fig:revised-pathway} are rearranged? If we assume that learners subconsciously keep track of the environments where they received direct evidence of gapping or resumption, which seems reasonable in light of growing evidence that they are sensitive to statistical patterns in the input (e.g.,~\cite{romberg.saffran2010}), then the principles presented in this chapter will guide the learner to the target grammar state regardless of how the steps are ordered. As we will see shortly, keeping track of direct evidence allows us to deal with cases where one relativization option needs to be overwritten by another. Take, for example, the acquisition pathway in Figure~\ref{fig:rearranged-pathway}, which represents the situation where the first piece of positive evidence of RCs the learner receives is for resumption at GEN.

\begin{figure}
\caption{L1 acquisition pathway for single-clause RCs from Figure \ref{fig:revised-pathway} with Steps~1 and 3 swapped}
\label{fig:rearranged-pathway}
\setstretch{1}
\raggedright
\begin{mybox}
{\bfseries Step 0: Initial state}
\par\vspace{2pt}
Assumed initial state has no resumption anywhere on the scale
\par\vspace{6pt}
\tikz{
    \draw[draw=white, line width=1pt] (0,-5pt) -- (0,5pt);
    \draw[line width=1pt] (0.4,0) -- (12.4,0);
    \draw[line width=1pt] (0.4,-5pt) -- (0.4,5pt);
    \draw[line width=1pt] (2.4,-5pt) -- (2.4,5pt);
    \draw[line width=1pt] (4.4,-5pt) -- (4.4,5pt);
    \draw[line width=1pt] (6.4,-5pt) -- (6.4,5pt);
    \draw[line width=1pt] (8.4,-5pt) -- (8.4,5pt);
    \draw[line width=1pt] (10.4,-5pt) -- (10.4,5pt);
    \draw[line width=1pt] (12.4,-5pt) -- (12.4,5pt);
    \path(1.4,-.3)node(){\strut\ttfamily SU};
    \path(3.4,-.3)node(){\strut\ttfamily DO};
    \path(5.4,-.3)node(){\strut\ttfamily IO};
    \path(7.4,-.3)node(){\strut\ttfamily OBL};
    \path(9.4,-.3)node(){\strut\ttfamily GEN};
    \path(11.4,-.3)node(){\strut\ttfamily OCOMP};
    \path(1.4,.3)node(){\strut\ttfamily 0};
    \path(3.4,.3)node(){\strut\ttfamily 0};
    \path(5.4,.3)node(){\strut\ttfamily 0};
    \path(7.4,.3)node(){\strut\ttfamily 0};
    \path(9.4,.3)node(){\strut\ttfamily 0};
    \path(11.4,.3)node(){\strut\ttfamily 0};
}\par
\vspace{3pt}
{\bfseries Step 1: Positive evidence of resumptive RCs at GEN}
\par\vspace{2pt}
Resumption added at GEN; Affordability demands that resumption be extended upward from GEN
\par\vspace{6pt}
\tikz{
    \draw[draw=white, line width=1pt] (0,-5pt) -- (0,5pt);
    \draw[line width=1pt] (0.4,0) -- (12.4,0);
    \draw[line width=1pt] (0.4,-5pt) -- (0.4,5pt);
    \draw[line width=1pt] (2.4,-5pt) -- (2.4,5pt);
    \draw[line width=1pt] (4.4,-5pt) -- (4.4,5pt);
    \draw[line width=1pt] (6.4,-5pt) -- (6.4,5pt);
    \draw[line width=1pt] (8.4,-5pt) -- (8.4,5pt);
    \draw[line width=1pt] (10.4,-5pt) -- (10.4,5pt);
    \draw[line width=1pt] (12.4,-5pt) -- (12.4,5pt);
    \path(1.4,-.3)node(){\strut\ttfamily SU};
    \path(3.4,-.3)node(){\strut\ttfamily DO};
    \path(5.4,-.3)node(){\strut\ttfamily IO};
    \path(7.4,-.3)node(){\strut\ttfamily OBL};
    \path(9.4,-.3)node(){\strut\ttfamily GEN};
    \path(11.4,-.3)node(){\strut\ttfamily OCOMP};
    \path(1.4,.3)node(){\strut\large\ttfamily +};
    \path(3.4,.3)node(){\strut\large\ttfamily +};
    \path(5.4,.3)node(){\strut\large\ttfamily +};
    \path(7.4,.3)node(){\strut\large\ttfamily +};
    \path(9.4,.3)node(){\strut\large\ttfamily [+]};
    \path(11.4,.3)node(){\strut\ttfamily 0};
}\par
\vspace{3pt}
{\bfseries Step 2: Positive evidence of gapped RCs at SU}
\par\vspace{2pt}
Gapping added at SU; Preemption demands that resumption be removed from SU
\par\vspace{6pt}
\tikz{
    \draw[draw=white, line width=1pt] (0,-5pt) -- (0,5pt);
    \draw[line width=1pt] (0.4,0) -- (12.4,0);
    \draw[line width=1pt] (0.4,-5pt) -- (0.4,5pt);
    \draw[line width=1pt] (2.4,-5pt) -- (2.4,5pt);
    \draw[line width=1pt] (4.4,-5pt) -- (4.4,5pt);
    \draw[line width=1pt] (6.4,-5pt) -- (6.4,5pt);
    \draw[line width=1pt] (8.4,-5pt) -- (8.4,5pt);
    \draw[line width=1pt] (10.4,-5pt) -- (10.4,5pt);
    \draw[line width=1pt] (12.4,-5pt) -- (12.4,5pt);
    \path(1.4,-.3)node(){\strut\ttfamily SU};
    \path(3.4,-.3)node(){\strut\ttfamily DO};
    \path(5.4,-.3)node(){\strut\ttfamily IO};
    \path(7.4,-.3)node(){\strut\ttfamily OBL};
    \path(9.4,-.3)node(){\strut\ttfamily GEN};
    \path(11.4,-.3)node(){\strut\ttfamily OCOMP};
    \path(1.4,.3)node(){\strut\large\ttfamily [−]};
    \path(3.4,.3)node(){\strut\large\ttfamily +};
    \path(5.4,.3)node(){\strut\large\ttfamily +};
    \path(7.4,.3)node(){\strut\large\ttfamily +};
    \path(9.4,.3)node(){\strut\large\ttfamily [+]};
    \path(11.4,.3)node(){\strut\ttfamily 0};
}\par
\vspace{3pt}
{\bfseries Step 3: Positive evidence of gapped RCs at IO }
\par\vspace{2pt}
Gapping added at IO; Continuity demands that gapping also be available at DO; Preemption demands that resumption be removed from DO and IO
\par\vspace{6pt}
\tikz{
    \draw[draw=white, line width=1pt] (0,-5pt) -- (0,5pt);
    \draw[line width=1pt] (0.4,0) -- (12.4,0);
    \draw[line width=1pt] (0.4,-5pt) -- (0.4,5pt);
    \draw[line width=1pt] (2.4,-5pt) -- (2.4,5pt);
    \draw[line width=1pt] (4.4,-5pt) -- (4.4,5pt);
    \draw[line width=1pt] (6.4,-5pt) -- (6.4,5pt);
    \draw[line width=1pt] (8.4,-5pt) -- (8.4,5pt);
    \draw[line width=1pt] (10.4,-5pt) -- (10.4,5pt);
    \draw[line width=1pt] (12.4,-5pt) -- (12.4,5pt);
    \path(1.4,-.3)node(){\strut\ttfamily SU};
    \path(3.4,-.3)node(){\strut\ttfamily DO};
    \path(5.4,-.3)node(){\strut\ttfamily IO};
    \path(7.4,-.3)node(){\strut\ttfamily OBL};
    \path(9.4,-.3)node(){\strut\ttfamily GEN};
    \path(11.4,-.3)node(){\strut\ttfamily OCOMP};
    \path(1.4,.3)node(){\strut\large\ttfamily [−]};
    \path(3.4,.3)node(){\strut\large\ttfamily −};
    \path(5.4,.3)node(){\strut\large\ttfamily [−]};
    \path(7.4,.3)node(){\strut\large\ttfamily +};
    \path(9.4,.3)node(){\strut\large\ttfamily [+]};
    \path(11.4,.3)node(){\strut\ttfamily 0};
}\par
\vspace{3pt}
{\bfseries Step 4: Positive evidence of resumptive RCs at DO}
\par\vspace{2pt}
Resumption added at DO; Continuity demands (a)~that gapping be available at SU through IO and (b)~that resumption be available at DO though GEN
\par\vspace{6pt}
\tikz{
    \draw[draw=white, line width=1pt] (0,-5pt) -- (0,5pt);
    \draw[line width=1pt] (0.4,0) -- (12.4,0);
    \draw[line width=1pt] (0.4,-5pt) -- (0.4,5pt);
    \draw[line width=1pt] (2.4,-5pt) -- (2.4,5pt);
    \draw[line width=1pt] (4.4,-5pt) -- (4.4,5pt);
    \draw[line width=1pt] (6.4,-5pt) -- (6.4,5pt);
    \draw[line width=1pt] (8.4,-5pt) -- (8.4,5pt);
    \draw[line width=1pt] (10.4,-5pt) -- (10.4,5pt);
    \draw[line width=1pt] (12.4,-5pt) -- (12.4,5pt);
    \path(1.4,-.3)node(){\strut\ttfamily SU};
    \path(3.4,-.3)node(){\strut\ttfamily DO};
    \path(5.4,-.3)node(){\strut\ttfamily IO};
    \path(7.4,-.3)node(){\strut\ttfamily OBL};
    \path(9.4,-.3)node(){\strut\ttfamily GEN};
    \path(11.4,-.3)node(){\strut\ttfamily OCOMP};
    \path(1.4,.3)node(){\strut\large\ttfamily [−]};
    \path(3.4,.3)node(){\strut\large\ttfamily -/[+]};
    \path(5.4,.3)node(){\strut\large\ttfamily [-]/+};
    \path(7.4,.3)node(){\strut\large\ttfamily +};
    \path(9.4,.3)node(){\strut\large\ttfamily [+]};
    \path(11.4,.3)node(){\strut\ttfamily 0};
}
\vspace{-6pt}
\end{mybox}
\vspace{-6pt}
\textit{Note}. 0 = unrelativizable; − = gapping; + = resumption; [ ] = direct evidence.\par
\setstretch{1.5}\vspace{-12pt}
\end{figure}

At Step~1, the learner responds to positive evidence of resumptive RCs at GEN by introducing resumption at that position. Affordability demands that resumption be extended upward from GEN to the top of the NPAH.

At Step~2, the learner responds to positive evidence of gapped RCs at SU by adding gapping at that position. Preemption then demands that resumption be removed from SU; in other words, resumption cannot be maintained at SU because only gapping is forced to be present there. Again, the only way for learners to establish whether an option is forced to be present (by either positive evidence or Continuity) is to track and retain information about the environments where there has been direct evidence for that option.

The same basic process is repeated at Step~3: The learner responds to positive evidence of gapped RCs at IO by adding gapping there. Continuity demands that gapping then be added at DO, and Preemption demands that resumption be removed from both DO and IO.

Finally, Step~4 proceeds just as it did in Figure~\ref{fig:revised-pathway}: The learner responds to positive evidence of resumptive RCs at DO by adding resumption there. Continuity demands (a)~that resumption be available at DO through GEN and (b)~that gapping be available at SU through IO. Thus, we have arrived again at the target grammar state with mandatory gapping at SU, optionality between gapping and resumption at DO through IO, mandatory resumption at OBL through GEN, and no relativization at OCOMP.

Figure \ref{fig:final-states} shows the L1 target states for English, Korean, and Mandarin in single-clause RCs according to \textcite{keenan.comrie1977}. These three target states can be accommodated easily in the framework I am proposing, each of them being simpler in some respects than the target grammar state from Figures~\ref{fig:revised-pathway} and \ref{fig:rearranged-pathway}. The diagram for English is the simplest one, with mandatory gapping across the entire scale. The diagram for Korean is slightly more complex, with mandatory gapping at SU through OBL, mandatory resumption at GEN, and no relativization at OCOMP. Finally, the diagram for Mandarin has mandatory gapping at SU (for single-clause RCs), optionality between gapping and resumption at DO, and mandatory resumption at IO through OCOMP. 

\begin{figure}
\caption{L1 target states for English, Korean, and Mandarin in single-clause RCs}
\label{fig:final-states}
\setstretch{1}
\raggedright
\begin{mybox}
{\bfseries English:}\par
\vspace{6pt}
\tikz{
    \draw[draw=white, line width=1pt] (0,-5pt) -- (0,5pt);
    \draw[line width=1pt] (0.4,0) -- (12.4,0);
    \draw[line width=1pt] (0.4,-5pt) -- (0.4,5pt);
    \draw[line width=1pt] (2.4,-5pt) -- (2.4,5pt);
    \draw[line width=1pt] (4.4,-5pt) -- (4.4,5pt);
    \draw[line width=1pt] (6.4,-5pt) -- (6.4,5pt);
    \draw[line width=1pt] (8.4,-5pt) -- (8.4,5pt);
    \draw[line width=1pt] (10.4,-5pt) -- (10.4,5pt);
    \draw[line width=1pt] (12.4,-5pt) -- (12.4,5pt);
    \path(1.4,-.3)node(){\strut\ttfamily SU};
    \path(3.4,-.3)node(){\strut\ttfamily DO};
    \path(5.4,-.3)node(){\strut\ttfamily IO};
    \path(7.4,-.3)node(){\strut\ttfamily OBL};
    \path(9.4,-.3)node(){\strut\ttfamily GEN};
    \path(11.4,-.3)node(){\strut\ttfamily OCOMP};
    \path(1.4,.3)node(){\strut\large\ttfamily −};
    \path(3.4,.3)node(){\strut\large\ttfamily −};
    \path(5.4,.3)node(){\strut\large\ttfamily −};
    \path(7.4,.3)node(){\strut\large\ttfamily −};
    \path(9.4,.3)node(){\strut\large\ttfamily −};
    \path(11.4,.3)node(){\strut\large\ttfamily −};
}\par
\vspace{3pt}
{\bfseries Korean:}\par
\vspace{6pt}
\tikz{
    \draw[draw=white, line width=1pt] (0,-5pt) -- (0,5pt);
    \draw[line width=1pt] (0.4,0) -- (12.4,0);
    \draw[line width=1pt] (0.4,-5pt) -- (0.4,5pt);
    \draw[line width=1pt] (2.4,-5pt) -- (2.4,5pt);
    \draw[line width=1pt] (4.4,-5pt) -- (4.4,5pt);
    \draw[line width=1pt] (6.4,-5pt) -- (6.4,5pt);
    \draw[line width=1pt] (8.4,-5pt) -- (8.4,5pt);
    \draw[line width=1pt] (10.4,-5pt) -- (10.4,5pt);
    \draw[line width=1pt] (12.4,-5pt) -- (12.4,5pt);
    \path(1.4,-.3)node(){\strut\ttfamily SU};
    \path(3.4,-.3)node(){\strut\ttfamily DO};
    \path(5.4,-.3)node(){\strut\ttfamily IO};
    \path(7.4,-.3)node(){\strut\ttfamily OBL};
    \path(9.4,-.3)node(){\strut\ttfamily GEN};
    \path(11.4,-.3)node(){\strut\ttfamily OCOMP};
    \path(1.4,.3)node(){\strut\large\ttfamily −};
    \path(3.4,.3)node(){\strut\large\ttfamily −};
    \path(5.4,.3)node(){\strut\large\ttfamily −};
    \path(7.4,.3)node(){\strut\large\ttfamily −};
    \path(9.4,.3)node(){\strut\large\ttfamily +};
    \path(11.4,.3)node(){\strut\ttfamily 0};
}\par
\vspace{3pt}
{\bfseries Mandarin:}\par
\vspace{6pt}
\tikz{
    \draw[draw=white, line width=1pt] (0,-5pt) -- (0,5pt);
    \draw[line width=1pt] (0.4,0) -- (12.4,0);
    \draw[line width=1pt] (0.4,-5pt) -- (0.4,5pt);
    \draw[line width=1pt] (2.4,-5pt) -- (2.4,5pt);
    \draw[line width=1pt] (4.4,-5pt) -- (4.4,5pt);
    \draw[line width=1pt] (6.4,-5pt) -- (6.4,5pt);
    \draw[line width=1pt] (8.4,-5pt) -- (8.4,5pt);
    \draw[line width=1pt] (10.4,-5pt) -- (10.4,5pt);
    \draw[line width=1pt] (12.4,-5pt) -- (12.4,5pt);
    \path(1.4,-.3)node(){\strut\ttfamily SU};
    \path(3.4,-.3)node(){\strut\ttfamily DO};
    \path(5.4,-.3)node(){\strut\ttfamily IO};
    \path(7.4,-.3)node(){\strut\ttfamily OBL};
    \path(9.4,-.3)node(){\strut\ttfamily GEN};
    \path(11.4,-.3)node(){\strut\ttfamily OCOMP};
    \path(1.4,.3)node(){\strut\large\ttfamily −};
    \path(3.4,.3)node(){\strut\large\ttfamily -/+};
    \path(5.4,.3)node(){\strut\large\ttfamily +};
    \path(7.4,.3)node(){\strut\large\ttfamily +};
    \path(9.4,.3)node(){\strut\large\ttfamily +};
    \path(11.4,.3)node(){\strut\large\ttfamily +};
}
\vspace{-6pt}
\end{mybox}
\vspace{-6pt}
\textit{Note}. 0 = unrelativizable; − = gapping; + = resumption.\par
\setstretch{1.5}\vspace{-12pt}
\end{figure} 

%========================================================%
\section{Extending the L1 acquisition pathway to the SLW hierarchy}\label{acq-slw-pathway}
%========================================================%

Although I have been focusing on the NPAH up to this point in the chapter, it is easy enough to apply the same approach to other processing hierarchies. Figure \ref{fig:slw-pathway} presents an L1 acquisition pathway for the SLW hierarchy in English ORCs. Note that whereas it is purely for the sake of exposition that I entertain the notion that every level on the NPAH constitutes a break point in difficulty/complexity, I consider each level on the SLW hierarchy to be a \textit{bona fide} break point: The `L' level is a break point because it requires the RC dependency to cross a clause boundary, and the `W' level is a break point because it requires the RC dependency to be extended into an embedded clause when there is another filler already waiting to be integrated (see~Section~\ref{acq-slw}).

\begin{figure}
\caption{L1 acquisition pathway for the SLW hierarchy in English ORCs}
\label{fig:slw-pathway}
\setstretch{1}
\raggedright
\begin{mybox}
{\bfseries Step 0: Initial state}
\par\vspace{2pt}
Assumed initial state has no resumption anywhere on the scale
\par\vspace{6pt}
\tikz{
    \draw[draw=white, line width=1pt] (0,-5pt) -- (0,5pt);
    \draw[line width=1pt] (0.4,0) -- (12.4,0);
    \draw[line width=1pt] (0.4,-5pt) -- (0.4,5pt);
    \draw[line width=1pt] (4.4,-5pt) -- (4.4,5pt);
    \draw[line width=1pt] (8.4,-5pt) -- (8.4,5pt);
    \draw[line width=1pt] (12.4,-5pt) -- (12.4,5pt);
    \path(2.4,-.3)node(){\strut\ttfamily\strut S};
    \path(6.4,-.3)node(){\strut\ttfamily\strut L};
    \path(10.4,-.3)node(){\strut\ttfamily\strut W};
    \path(2.4,.3)node(){\strut\ttfamily 0};
    \path(6.4,.3)node(){\strut\ttfamily 0};
    \path(10.4,.3)node(){\strut\ttfamily 0};
}\par
\vspace{3pt}
{\bfseries Step 1: Positive evidence of gapping in short-distance ORCs} 
\par\vspace{2pt}
Gapping added at S
\par\vspace{6pt}
\tikz{
    \draw[draw=white, line width=1pt] (0,-5pt) -- (0,5pt);
    \draw[line width=1pt] (0.4,0) -- (12.4,0);
    \draw[line width=1pt] (0.4,-5pt) -- (0.4,5pt);
    \draw[line width=1pt] (4.4,-5pt) -- (4.4,5pt);
    \draw[line width=1pt] (8.4,-5pt) -- (8.4,5pt);
    \draw[line width=1pt] (12.4,-5pt) -- (12.4,5pt);
    \path(2.4,-.3)node(){\strut\ttfamily\strut S};
    \path(6.4,-.3)node(){\strut\ttfamily\strut L};
    \path(10.4,-.3)node(){\strut\ttfamily\strut W};
    \path(2.4,.3)node(){\strut\large\ttfamily [−]};
    \path(6.4,.3)node(){\strut\ttfamily 0};
    \path(10.4,.3)node(){\strut\ttfamily 0};
}\par
\vspace{3pt}
{\bfseries Step 2: Positive evidence of gapping in long-distance ORCs} 
\par\vspace{2pt}
Gapping added at L
\par\vspace{6pt}
\tikz{
    \draw[draw=white, line width=1pt] (0,-5pt) -- (0,5pt);
    \draw[line width=1pt] (0.4,0) -- (12.4,0);
    \draw[line width=1pt] (0.4,-5pt) -- (0.4,5pt);
    \draw[line width=1pt] (4.4,-5pt) -- (4.4,5pt);
    \draw[line width=1pt] (8.4,-5pt) -- (8.4,5pt);
    \draw[line width=1pt] (12.4,-5pt) -- (12.4,5pt);
    \path(2.4,-.3)node(){\strut\ttfamily\strut S};
    \path(6.4,-.3)node(){\strut\ttfamily\strut L};
    \path(10.4,-.3)node(){\strut\ttfamily\strut W};
    \path(2.4,.3)node(){\strut\large\ttfamily [−]};
    \path(6.4,.3)node(){\strut\large\ttfamily [−]};
    \path(10.4,.3)node(){\strut\ttfamily 0};
}
\vspace{-6pt}
\end{mybox}
\vspace{-6pt}
\textit{Note}. 0 = unrelativizable; − = gapping; + = resumption; [ ] = direct evidence.\par
\setstretch{1.5}\vspace{-12pt}
\end{figure}

Step~0 represents the initial state with no relativization permitted anywhere on the hierarchy. In Step~1, the learner responds to positive evidence of gapping in short-distance ORCs by adding gapping at S. In Step~2, the learner responds to positive evidence of gapping in long-distance ORCs by adding gapping at L.\footnotemark{} This yields the target grammar state for English ORCs, with mandatory gapping at S through L and no relativization at W.

\footnotetext{One might wonder what happens with cases of serial embedding, since any language that allows a first level of embedding for COMP-clauses allows unbounded levels of such embedding. My position on this is that once the learner starts allowing filler reactivation in COMP-clauses (see~Section~\ref{pro-routine}), there is nothing preventing a filler from being reactivated multiple times in the same sentence; although each embedding surely incurs a processing cost, it does not constitute a break point under Gatekeeping.}

Figure \ref{fig:slw-orc} shows the L1 target states for English, Korean, and Mandarin. As we have already established, English has mandatory gapping at S through L and no relativization at W. The L1 target states for Korean and Mandarin are based on the judgments of my language consultants. Korean is the same as English except that it is impossible in that language to replicate the situation where filler needs to be reactivated in an embedded question introduced by a fronted \textit{wh}-phrase because (a)~Korean has gap-filler RCs, so no filler has yet been encountered when an RC dependency crosses a clause boundary, and (b)~Korean lacks overt \textit{wh}-movement in questions, so embedded questions are not introduced by fillers (see~Section~\ref{gap-lw}). Thus, the `W' level is not applicable (I~use `N/A' to indicate that the `W' level on the hierarchy is absent), and English \textit{wh}-island ORCs translated into Korean should behave the same as regular long-distance ORCs. Finally, Mandarin is the same as Korean except that it optionally allows resumption at S and L. I use `N/A' to indicate that the `W' level on the hierarchy is absent from Korean and Mandarin.

\begin{figure}
\caption{L1 target states for English, Korean, and Mandarin on the SLW hierarchy for ORCs}
\label{fig:slw-orc}
\setstretch{1}
\raggedright
\begin{mybox}
{\bfseries English:}\par
\vspace{6pt}
\tikz{
    \draw[draw=white, line width=1pt] (0,-5pt) -- (0,5pt);
    \draw[line width=1pt] (0.4,0) -- (12.4,0);
    \draw[line width=1pt] (0.4,-5pt) -- (0.4,5pt);
    \draw[line width=1pt] (4.4,-5pt) -- (4.4,5pt);
    \draw[line width=1pt] (8.4,-5pt) -- (8.4,5pt);
    \draw[line width=1pt] (12.4,-5pt) -- (12.4,5pt);
    \path(2.4,-.3)node(){\strut\ttfamily\strut S};
    \path(6.4,-.3)node(){\strut\ttfamily\strut L};
    \path(10.4,-.3)node(){\strut\ttfamily\strut W};
    \path(2.4,.3)node(){\strut\large\ttfamily −};
    \path(6.4,.3)node(){\strut\large\ttfamily −};
    \path(10.4,.3)node(){\strut\ttfamily 0};
}\par
\vspace{3pt}
{\bfseries Korean:}\par
\vspace{6pt}
\tikz{
    \draw[draw=white, line width=1pt] (0,-5pt) -- (0,5pt);
    \draw[line width=1pt] (0.4,0) -- (12.4,0);
    \draw[line width=1pt] (0.4,-5pt) -- (0.4,5pt);
    \draw[line width=1pt] (4.4,-5pt) -- (4.4,5pt);
    \draw[line width=1pt] (8.4,-5pt) -- (8.4,5pt);
    \draw[line width=1pt] (12.4,-5pt) -- (12.4,5pt);
    \path(2.4,-.3)node(){\strut\ttfamily\strut S};
    \path(6.4,-.3)node(){\strut\ttfamily\strut L};
    \path(10.4,-.3)node(){\strut\ttfamily\strut W};
    \path(2.4,.3)node(){\strut\large\ttfamily −};
    \path(6.4,.3)node(){\strut\large\ttfamily −};
    \path(10.4,.3)node(){\strut\ttfamily N/A};
}\par
\vspace{3pt}
{\bfseries Mandarin:}\par
\vspace{6pt}
\tikz{
    \draw[draw=white, line width=1pt] (0,-5pt) -- (0,5pt);
    \draw[line width=1pt] (0.4,0) -- (12.4,0);
    \draw[line width=1pt] (0.4,-5pt) -- (0.4,5pt);
    \draw[line width=1pt] (4.4,-5pt) -- (4.4,5pt);
    \draw[line width=1pt] (8.4,-5pt) -- (8.4,5pt);
    \draw[line width=1pt] (12.4,-5pt) -- (12.4,5pt);
    \path(2.4,-.3)node(){\strut\ttfamily\strut S};
    \path(6.4,-.3)node(){\strut\ttfamily\strut L};
    \path(10.4,-.3)node(){\strut\ttfamily\strut W};
    \path(2.4,.3)node(){\strut\large\ttfamily -/+};
    \path(6.4,.3)node(){\strut\large\ttfamily -/+};
    \path(10.4,.3)node(){\strut\ttfamily N/A};
}
\vspace{-6pt}
\end{mybox}
\vspace{-6pt}
\textit{Note}. 0 = unrelativizable; − = gapping; + = resumption; N/A = not applicable.\par
\setstretch{1.5}\vspace{-12pt}
\end{figure}

Figure \ref{fig:slw-src} shows the SRC facts for purposes of comparison. The target states for English and Korean are the same here as they were for ORCs. The one difference is in the diagram for Mandarin, which has mandatory gapping instead of optionality between gapping and resumption at S, reflecting the fact that Mandarin requires gapping in short-distance SRCs.

\begin{figure}
\caption{L1 target states for English, Korean, and Mandarin on the SLW hierarchy for SRCs}
\label{fig:slw-src}
\setstretch{1}
\raggedright
\begin{mybox}
{\bfseries English:}\par
\vspace{6pt}
\tikz{
    \draw[draw=white, line width=1pt] (0,-5pt) -- (0,5pt);
    \draw[line width=1pt] (0.4,0) -- (12.4,0);
    \draw[line width=1pt] (0.4,-5pt) -- (0.4,5pt);
    \draw[line width=1pt] (4.4,-5pt) -- (4.4,5pt);
    \draw[line width=1pt] (8.4,-5pt) -- (8.4,5pt);
    \draw[line width=1pt] (12.4,-5pt) -- (12.4,5pt);
    \path(2.4,-.3)node(){\strut\ttfamily\strut S};
    \path(6.4,-.3)node(){\strut\ttfamily\strut L};
    \path(10.4,-.3)node(){\strut\ttfamily\strut W};
    \path(2.4,.3)node(){\strut\large\ttfamily −};
    \path(6.4,.3)node(){\strut\large\ttfamily −};
    \path(10.4,.3)node(){\strut\ttfamily 0};
}\par
\vspace{3pt}
{\bfseries Korean:}\par
\vspace{6pt}
\tikz{
    \draw[draw=white, line width=1pt] (0,-5pt) -- (0,5pt);
    \draw[line width=1pt] (0.4,0) -- (12.4,0);
    \draw[line width=1pt] (0.4,-5pt) -- (0.4,5pt);
    \draw[line width=1pt] (4.4,-5pt) -- (4.4,5pt);
    \draw[line width=1pt] (8.4,-5pt) -- (8.4,5pt);
    \draw[line width=1pt] (12.4,-5pt) -- (12.4,5pt);
    \path(2.4,-.3)node(){\strut\ttfamily\strut S};
    \path(6.4,-.3)node(){\strut\ttfamily\strut L};
    \path(10.4,-.3)node(){\strut\ttfamily\strut W};
    \path(2.4,.3)node(){\strut\large\ttfamily −};
    \path(6.4,.3)node(){\strut\large\ttfamily −};
    \path(10.4,.3)node(){\strut\ttfamily N/A};
}\par
\vspace{3pt}
{\bfseries Mandarin:}\par
\vspace{6pt}
\tikz{
    \draw[draw=white, line width=1pt] (0,-5pt) -- (0,5pt);
    \draw[line width=1pt] (0.4,0) -- (12.4,0);
    \draw[line width=1pt] (0.4,-5pt) -- (0.4,5pt);
    \draw[line width=1pt] (4.4,-5pt) -- (4.4,5pt);
    \draw[line width=1pt] (8.4,-5pt) -- (8.4,5pt);
    \draw[line width=1pt] (12.4,-5pt) -- (12.4,5pt);
    \path(2.4,-.3)node(){\strut\ttfamily\strut S};
    \path(6.4,-.3)node(){\strut\ttfamily\strut L};
    \path(10.4,-.3)node(){\strut\ttfamily\strut W};
    \path(2.4,.3)node(){\strut\large\ttfamily −};
    \path(6.4,.3)node(){\strut\large\ttfamily -/+};
    \path(10.4,.3)node(){\strut\ttfamily N/A};
}
\vspace{-6pt}
\end{mybox}
\vspace{-6pt}
\textit{Note}. 0 = unrelativizable; − = gapping; + = resumption; N/A = not applicable.\par
\setstretch{1.5}\vspace{-12pt}
\end{figure}

%========================================================%
\section{Reconsidering the L2 initial state}\label{acq-initial-state}
%========================================================%

As we have already seen, the claims that \textcite{hyltenstam1984,hyltenstam1990} and \textcite{eckman2004} made about the L2 initial state start to show cracks when subjected to careful scrutiny. For starters, treating resumption as the default option for relativization in L1/L2 acquisition is problematic for at least two reasons. First, there is strong experimental evidence indicating that gapping is easier than resumption at the top of the processing hierarchies considered in this dissertation. Second, such an approach fails to accommodate languages that lack relativization altogether in certain domains. 

These first two issues can be dealt with by altering our view of the L1 initial state so that it allows for no relativization anywhere on the hierarchy, in line with Conservatism. Even after making this change, however, claiming that the L2 initial state is always characterized by a more or less full reset to the assumed L1 initial state (\cite{eckman2004}, p.~106ff; \cite{hyltenstam1984}, p.~43) is still problematic because it fails to account for L1 influence of the kind reported by \textcite{hyltenstam1984,hyltenstam1990}, \textcite{hawkins.chan1997}, and others in this line of research. Another possibility is that L2ers start by `regressing' down the hierarchy only part-way towards the L1 initial state, as \textcite{hyltenstam1984} suggested might happen (p.~55), but no proposal has been put forward to my knowledge which explains either how far L2ers can be expected to regress or what prevents them from going further. 

Another possibility is offered by the Full Transfer/Full Access (FT/FA) model (\cite{schwartz.sprouse1994,schwartz.sprouse1996}): The L2 initial state is identical to the current L1 state at the onset of L2 acquisition. One advantage that this approach has over Hyltenstam's partial-reset approach is that it gives clear-cut predictions about what the initial state should look like for each L2er, which makes it more readily falsifiable. In the present dissertation, I test L2ers in both their L1 (Korean or Mandarin) and the target language (English) so that L1 transfer effects can be identified; a pattern of results indicating that the L1-Mandarin L2ers had transferred grammatical resumption to their interlanguages would be consistent with the full-transfer approach to the L2 initial state (and perhaps the partial-reset approach) but not with a full-reset approach where learners revert to an L1 initial state that has no relativization across all environments.

Assuming for a moment that the full-transfer approach is on the right track, let us consider now what the resulting processing-guided L2 acquisition pathways might look like for several language pairings we have been discussing in this chapter. The easiest cases to deal with are those where gapping or resumption merely needs to be extended into a domain where no relativization is permitted in the L1. Take, for example, L1-Finnish L2ers of Swedish. As Figure \ref{fig:challenge-swedish} illustrates, all that is required in such cases is for the learner to extend gapping down through OCOMP, reflecting a grammar state with mandatory gapping across the entire scale.

\begin{figure}
\caption{Acquisitional task facing L1-Finnish L2ers of Swedish in single-clause RCs with full transfer}
\label{fig:challenge-swedish}
\setstretch{1}
\raggedright
\begin{mybox}
{\bfseries L2 initial state:}\par
\vspace{6pt}
\tikz{
    \draw[draw=white, line width=1pt] (0,-5pt) -- (0,5pt);
    \draw[line width=1pt] (0.4,0) -- (12.4,0);
    \draw[line width=1pt] (0.4,-5pt) -- (0.4,5pt);
    \draw[line width=1pt] (2.4,-5pt) -- (2.4,5pt);
    \draw[line width=1pt] (4.4,-5pt) -- (4.4,5pt);
    \draw[line width=1pt] (6.4,-5pt) -- (6.4,5pt);
    \draw[line width=1pt] (8.4,-5pt) -- (8.4,5pt);
    \draw[line width=1pt] (10.4,-5pt) -- (10.4,5pt);
    \draw[line width=1pt] (12.4,-5pt) -- (12.4,5pt);
    \path(1.4,-.3)node(){\strut\ttfamily SU};
    \path(3.4,-.3)node(){\strut\ttfamily DO};
    \path(5.4,-.3)node(){\strut\ttfamily IO};
    \path(7.4,-.3)node(){\strut\ttfamily OBL};
    \path(9.4,-.3)node(){\strut\ttfamily GEN};
    \path(11.4,-.3)node(){\strut\ttfamily OCOMP};
    \path(1.4,.3)node(){\strut\large\ttfamily −};
    \path(3.4,.3)node(){\strut\large\ttfamily −};
    \path(5.4,.3)node(){\strut\large\ttfamily −};
    \path(7.4,.3)node(){\strut\large\ttfamily −};
    \path(9.4,.3)node(){\strut\large\ttfamily −};
    \path(11.4,.3)node(){\strut\ttfamily 0};
}\par
\vspace{3pt}
{\bfseries L2 target state:}\par
\vspace{6pt}
\tikz{
    \draw[draw=white, line width=1pt] (0,-5pt) -- (0,5pt);
    \draw[line width=1pt] (0.4,0) -- (12.4,0);
    \draw[line width=1pt] (0.4,-5pt) -- (0.4,5pt);
    \draw[line width=1pt] (2.4,-5pt) -- (2.4,5pt);
    \draw[line width=1pt] (4.4,-5pt) -- (4.4,5pt);
    \draw[line width=1pt] (6.4,-5pt) -- (6.4,5pt);
    \draw[line width=1pt] (8.4,-5pt) -- (8.4,5pt);
    \draw[line width=1pt] (10.4,-5pt) -- (10.4,5pt);
    \draw[line width=1pt] (12.4,-5pt) -- (12.4,5pt);
    \path(1.4,-.3)node(){\strut\ttfamily SU};
    \path(3.4,-.3)node(){\strut\ttfamily DO};
    \path(5.4,-.3)node(){\strut\ttfamily IO};
    \path(7.4,-.3)node(){\strut\ttfamily OBL};
    \path(9.4,-.3)node(){\strut\ttfamily GEN};
    \path(11.4,-.3)node(){\strut\ttfamily OCOMP};
    \path(1.4,.3)node(){\strut\large\ttfamily −};
    \path(3.4,.3)node(){\strut\large\ttfamily −};
    \path(5.4,.3)node(){\strut\large\ttfamily −};
    \path(7.4,.3)node(){\strut\large\ttfamily −};
    \path(9.4,.3)node(){\strut\large\ttfamily −};
    \path(11.4,.3)node(){\strut\large\ttfamily −};
}
\vspace{-6pt}
\end{mybox}
\vspace{-6pt}
\textit{Note}. 0 = unrelativizable; − = gapping; + = resumption.\par
\setstretch{1.5}\vspace{-12pt}
\end{figure}

\noindent Proceeding from the L2 initial state to the target state requires only positive evidence of gapping at OCOMP, as shown in Figure \ref{fig:restructure-swedish}.

\begin{figure}
\caption{L2 acquisition pathway for L1-Finnish L2ers of Swedish in single-clause RCs}
\label{fig:restructure-swedish}
\setstretch{1}
\raggedright
\begin{mybox}
{\bfseries Step 0: L2 initial state with full transfer}
\par\vspace{2pt}
Assumed initial state has gapping at SU through GEN and no relativization at OCOMP
\par\vspace{6pt}
\tikz{
    \draw[draw=white, line width=1pt] (0,-5pt) -- (0,5pt);
    \draw[line width=1pt] (0.4,0) -- (12.4,0);
    \draw[line width=1pt] (0.4,-5pt) -- (0.4,5pt);
    \draw[line width=1pt] (2.4,-5pt) -- (2.4,5pt);
    \draw[line width=1pt] (4.4,-5pt) -- (4.4,5pt);
    \draw[line width=1pt] (6.4,-5pt) -- (6.4,5pt);
    \draw[line width=1pt] (8.4,-5pt) -- (8.4,5pt);
    \draw[line width=1pt] (10.4,-5pt) -- (10.4,5pt);
    \draw[line width=1pt] (12.4,-5pt) -- (12.4,5pt);
    \path(1.4,-.3)node(){\strut\ttfamily SU};
    \path(3.4,-.3)node(){\strut\ttfamily DO};
    \path(5.4,-.3)node(){\strut\ttfamily IO};
    \path(7.4,-.3)node(){\strut\ttfamily OBL};
    \path(9.4,-.3)node(){\strut\ttfamily GEN};
    \path(11.4,-.3)node(){\strut\ttfamily OCOMP};
    \path(1.4,.3)node(){\strut\large\ttfamily −};
    \path(3.4,.3)node(){\strut\large\ttfamily −};
    \path(5.4,.3)node(){\strut\large\ttfamily −};
    \path(7.4,.3)node(){\strut\large\ttfamily −};
    \path(9.4,.3)node(){\strut\large\ttfamily −};
    \path(11.4,.3)node(){\strut\ttfamily 0};
}\par
\vspace{3pt}
{\bfseries Step 1: Positive evidence of gapped RCs at OCOMP} 
\par\vspace{2pt}
Gapping added at OCOMP
\par\vspace{6pt}
\tikz{
    \draw[draw=white, line width=1pt] (0,-5pt) -- (0,5pt);
    \draw[line width=1pt] (0.4,0) -- (12.4,0);
    \draw[line width=1pt] (0.4,-5pt) -- (0.4,5pt);
    \draw[line width=1pt] (2.4,-5pt) -- (2.4,5pt);
    \draw[line width=1pt] (4.4,-5pt) -- (4.4,5pt);
    \draw[line width=1pt] (6.4,-5pt) -- (6.4,5pt);
    \draw[line width=1pt] (8.4,-5pt) -- (8.4,5pt);
    \draw[line width=1pt] (10.4,-5pt) -- (10.4,5pt);
    \draw[line width=1pt] (12.4,-5pt) -- (12.4,5pt);
    \path(1.4,-.3)node(){\strut\ttfamily SU};
    \path(3.4,-.3)node(){\strut\ttfamily DO};
    \path(5.4,-.3)node(){\strut\ttfamily IO};
    \path(7.4,-.3)node(){\strut\ttfamily OBL};
    \path(9.4,-.3)node(){\strut\ttfamily GEN};
    \path(11.4,-.3)node(){\strut\ttfamily OCOMP};
    \path(1.4,.3)node(){\strut\large\ttfamily −};
    \path(3.4,.3)node(){\strut\large\ttfamily −};
    \path(5.4,.3)node(){\strut\large\ttfamily −};
    \path(7.4,.3)node(){\strut\large\ttfamily −};
    \path(9.4,.3)node(){\strut\large\ttfamily −};
    \path(11.4,.3)node(){\strut\large\ttfamily [−]};
}
\vspace{-6pt}
\end{mybox}
\vspace{-6pt}
\textit{Note}. 0 = unrelativizable; − = gapping.\par
\setstretch{1.5}\vspace{-12pt}
\end{figure}

So far, so good. But arriving at the target state could be more challenging when gapping or resumption needs to be removed from domains where it is permitted in the L1. Consider the case of an L1-Mandarin L2er of English. Figure \ref{fig:challenge-english} shows that in addition to extending gapping down through OCOMP, the L2er must also purge resumption from DO through OCOMP. 

\begin{figure}
\caption{Acquisitional task facing L1-Mandarin L2ers of English in single-clause RCs with full transfer}
\label{fig:challenge-english}
\setstretch{1}
\raggedright
\begin{mybox}
{\bfseries L2 initial state:}\par
\vspace{6pt}
\tikz{
    \draw[draw=white, line width=1pt] (0,-5pt) -- (0,5pt);
    \draw[line width=1pt] (0.4,0) -- (12.4,0);
    \draw[line width=1pt] (0.4,-5pt) -- (0.4,5pt);
    \draw[line width=1pt] (2.4,-5pt) -- (2.4,5pt);
    \draw[line width=1pt] (4.4,-5pt) -- (4.4,5pt);
    \draw[line width=1pt] (6.4,-5pt) -- (6.4,5pt);
    \draw[line width=1pt] (8.4,-5pt) -- (8.4,5pt);
    \draw[line width=1pt] (10.4,-5pt) -- (10.4,5pt);
    \draw[line width=1pt] (12.4,-5pt) -- (12.4,5pt);
    \path(1.4,-.3)node(){\strut\ttfamily SU};
    \path(3.4,-.3)node(){\strut\ttfamily DO};
    \path(5.4,-.3)node(){\strut\ttfamily IO};
    \path(7.4,-.3)node(){\strut\ttfamily OBL};
    \path(9.4,-.3)node(){\strut\ttfamily GEN};
    \path(11.4,-.3)node(){\strut\ttfamily OCOMP};
    \path(1.4,.3)node(){\strut\large\ttfamily −};
    \path(3.4,.3)node(){\strut\large\ttfamily -/+};
    \path(5.4,.3)node(){\strut\large\ttfamily +};
    \path(7.4,.3)node(){\strut\large\ttfamily +};
    \path(9.4,.3)node(){\strut\large\ttfamily +};
    \path(11.4,.3)node(){\strut\large\ttfamily +};
}\par
\vspace{3pt}
{\bfseries L2 target state:}\par
\vspace{6pt}
\tikz{
    \draw[draw=white, line width=1pt] (0,-5pt) -- (0,5pt);
    \draw[line width=1pt] (0.4,0) -- (12.4,0);
    \draw[line width=1pt] (0.4,-5pt) -- (0.4,5pt);
    \draw[line width=1pt] (2.4,-5pt) -- (2.4,5pt);
    \draw[line width=1pt] (4.4,-5pt) -- (4.4,5pt);
    \draw[line width=1pt] (6.4,-5pt) -- (6.4,5pt);
    \draw[line width=1pt] (8.4,-5pt) -- (8.4,5pt);
    \draw[line width=1pt] (10.4,-5pt) -- (10.4,5pt);
    \draw[line width=1pt] (12.4,-5pt) -- (12.4,5pt);
    \path(1.4,-.3)node(){\strut\ttfamily SU};
    \path(3.4,-.3)node(){\strut\ttfamily DO};
    \path(5.4,-.3)node(){\strut\ttfamily IO};
    \path(7.4,-.3)node(){\strut\ttfamily OBL};
    \path(9.4,-.3)node(){\strut\ttfamily GEN};
    \path(11.4,-.3)node(){\strut\ttfamily OCOMP};
    \path(1.4,.3)node(){\strut\large\ttfamily −};
    \path(3.4,.3)node(){\strut\large\ttfamily −};
    \path(5.4,.3)node(){\strut\large\ttfamily −};
    \path(7.4,.3)node(){\strut\large\ttfamily −};
    \path(9.4,.3)node(){\strut\large\ttfamily −};
    \path(11.4,.3)node(){\strut\large\ttfamily −};
}
\vspace{-6pt} 
\end{mybox}
\vspace{-6pt} 
\textit{Note}. 0 = unrelativizable; − = gapping; + = resumption.\par
\setstretch{1.5}\vspace{-12pt}
\end{figure}

How might the L2er remove gapping and resumption from domains where they are permitted in the L1 but not in the target language? Some L2 researchers have suggested that explicit negative evidence (i.e.,~corrective feedback) is the solution to problems of this kind (\cite{bley-vroman1990,carroll.swain1993,ellis.etal2006,white1987,white1991}). Others have suggested that L2ers make use of indirect negative evidence (i.e.,~the absence in the input of positive evidence where it would be expected to occur) to determine what is disallowed in the target language (\cite{archibald1993,gass1988,plough1994}). However, there is disagreement among acquisitionists as to whether negative evidence, be it explicit or indirect, contributes to shaping L2 competence (\cite{gass2003}, p.~226; \cite{li.vuono2019}, p.~95; \cite{sylte2021}, p.~1). \textcite{schwartz1993}, for instance, argues that explicit negative evidence can only influence metalinguistic knowledge and has no effect on underlying linguistic competence (p.~150ff; see~also \cite{schwartz1986,schwartz1987}; \cite{schwartz.gubala-ryzak1992}; \cite{truscott1998}).

Whether the approach to RC acquisition I have presented in this chapter allows for the purging of values (i.e.,~gapping and resumption) transferred from the L1 based on target-language positive evidence alone hinges on whether the `supported by direct evidence' designation (represented by square brackets in the figures) carries over as part of L1 transfer. First, let us consider Figure~\ref{fig:additive}, which shows how acquisition of English single-clause RCs unfolds for L1-Mandarin L2ers if transferred values retain the `supported by direct evidence' designation. I call this the `additive' approach because it allows only for the addition of new values and not for the deletion of transferred values.

% \footnotetext{This can also be framed as a question of whether information about statistical patterns in the input transfers from the L1 to the L2.}

\begin{figure}
\caption{The additive approach to restructuring for L1-Mandarin L2ers of English in single-clause RCs}
\label{fig:additive}
\setstretch{1}
\raggedright
\begin{mybox}
{\bfseries Step 0: L2 initial state with full transfer}
\par\vspace{2pt}
Assumed initial state has gapping at SU, optionality at DO, and resumption at IO through OCOMP
\par\vspace{6pt}
\tikz{
    \draw[draw=white, line width=1pt] (0,-5pt) -- (0,5pt);
    \draw[line width=1pt] (0.4,0) -- (12.4,0);
    \draw[line width=1pt] (0.4,-5pt) -- (0.4,5pt);
    \draw[line width=1pt] (2.4,-5pt) -- (2.4,5pt);
    \draw[line width=1pt] (4.4,-5pt) -- (4.4,5pt);
    \draw[line width=1pt] (6.4,-5pt) -- (6.4,5pt);
    \draw[line width=1pt] (8.4,-5pt) -- (8.4,5pt);
    \draw[line width=1pt] (10.4,-5pt) -- (10.4,5pt);
    \draw[line width=1pt] (12.4,-5pt) -- (12.4,5pt);
    \path(1.4,-.3)node(){\strut\ttfamily SU};
    \path(3.4,-.3)node(){\strut\ttfamily DO};
    \path(5.4,-.3)node(){\strut\ttfamily IO};
    \path(7.4,-.3)node(){\strut\ttfamily OBL};
    \path(9.4,-.3)node(){\strut\ttfamily GEN};
    \path(11.4,-.3)node(){\strut\ttfamily OCOMP};
    \path(1.4,.3)node(){\strut\large\ttfamily [−]};
    \path(3.4,.3)node(){\strut\large\ttfamily [-]/[+]};
    \path(5.4,.3)node(){\strut\large\ttfamily [+]};
    \path(7.4,.3)node(){\strut\large\ttfamily [+]};
    \path(9.4,.3)node(){\strut\large\ttfamily [+]};
    \path(11.4,.3)node(){\strut\large\ttfamily [+]};
}\par
\vspace{3pt}
{\bfseries Step 1: Positive evidence of gapped RCs at SU}
\par\vspace{2pt}
No change
\par\vspace{6pt}
\tikz{
    \draw[draw=white, line width=1pt] (0,-5pt) -- (0,5pt);
    \draw[line width=1pt] (0.4,0) -- (12.4,0);
    \draw[line width=1pt] (0.4,-5pt) -- (0.4,5pt);
    \draw[line width=1pt] (2.4,-5pt) -- (2.4,5pt);
    \draw[line width=1pt] (4.4,-5pt) -- (4.4,5pt);
    \draw[line width=1pt] (6.4,-5pt) -- (6.4,5pt);
    \draw[line width=1pt] (8.4,-5pt) -- (8.4,5pt);
    \draw[line width=1pt] (10.4,-5pt) -- (10.4,5pt);
    \draw[line width=1pt] (12.4,-5pt) -- (12.4,5pt);
    \path(1.4,-.3)node(){\strut\ttfamily SU};
    \path(3.4,-.3)node(){\strut\ttfamily DO};
    \path(5.4,-.3)node(){\strut\ttfamily IO};
    \path(7.4,-.3)node(){\strut\ttfamily OBL};
    \path(9.4,-.3)node(){\strut\ttfamily GEN};
    \path(11.4,-.3)node(){\strut\ttfamily OCOMP};
    \path(1.4,.3)node(){\strut\large\ttfamily [−]};
    \path(3.4,.3)node(){\strut\large\ttfamily [-]/[+]};
    \path(5.4,.3)node(){\strut\large\ttfamily [+]};
    \path(7.4,.3)node(){\strut\large\ttfamily [+]};
    \path(9.4,.3)node(){\strut\large\ttfamily [+]};
    \path(11.4,.3)node(){\strut\large\ttfamily [+]};
}\par\vspace{3pt}
{\bfseries Step 2: Positive evidence of gapped RCs at IO}
\par\vspace{2pt}
Gapping added at IO
\par\vspace{6pt}
\tikz{
    \draw[draw=white, line width=1pt] (0,-5pt) -- (0,5pt);
    \draw[line width=1pt] (0.4,0) -- (12.4,0);
    \draw[line width=1pt] (0.4,-5pt) -- (0.4,5pt);
    \draw[line width=1pt] (2.4,-5pt) -- (2.4,5pt);
    \draw[line width=1pt] (4.4,-5pt) -- (4.4,5pt);
    \draw[line width=1pt] (6.4,-5pt) -- (6.4,5pt);
    \draw[line width=1pt] (8.4,-5pt) -- (8.4,5pt);
    \draw[line width=1pt] (10.4,-5pt) -- (10.4,5pt);
    \draw[line width=1pt] (12.4,-5pt) -- (12.4,5pt);
    \path(1.4,-.3)node(){\strut\ttfamily SU};
    \path(3.4,-.3)node(){\strut\ttfamily DO};
    \path(5.4,-.3)node(){\strut\ttfamily IO};
    \path(7.4,-.3)node(){\strut\ttfamily OBL};
    \path(9.4,-.3)node(){\strut\ttfamily GEN};
    \path(11.4,-.3)node(){\strut\ttfamily OCOMP};
    \path(1.4,.3)node(){\strut\large\ttfamily [−]};
    \path(3.4,.3)node(){\strut\large\ttfamily [-]/[+]};
    \path(5.4,.3)node(){\strut\large\ttfamily [-]/[+]};
    \path(7.4,.3)node(){\strut\large\ttfamily [+]};
    \path(9.4,.3)node(){\strut\large\ttfamily [+]};
    \path(11.4,.3)node(){\strut\large\ttfamily [+]};
}\par\vspace{3pt}
{\bfseries Step 3: Positive evidence of gapped RCs at OCOMP}
\par\vspace{2pt}
Gapping added at OCOMP; Continuity demands that gapping also be added at OBL and GEN
\par\vspace{6pt}
\tikz{
    \draw[draw=white, line width=1pt] (0,-5pt) -- (0,5pt);
    \draw[line width=1pt] (0.4,0) -- (12.4,0);
    \draw[line width=1pt] (0.4,-5pt) -- (0.4,5pt);
    \draw[line width=1pt] (2.4,-5pt) -- (2.4,5pt);
    \draw[line width=1pt] (4.4,-5pt) -- (4.4,5pt);
    \draw[line width=1pt] (6.4,-5pt) -- (6.4,5pt);
    \draw[line width=1pt] (8.4,-5pt) -- (8.4,5pt);
    \draw[line width=1pt] (10.4,-5pt) -- (10.4,5pt);
    \draw[line width=1pt] (12.4,-5pt) -- (12.4,5pt);
    \path(1.4,-.3)node(){\strut\ttfamily SU};
    \path(3.4,-.3)node(){\strut\ttfamily DO};
    \path(5.4,-.3)node(){\strut\ttfamily IO};
    \path(7.4,-.3)node(){\strut\ttfamily OBL};
    \path(9.4,-.3)node(){\strut\ttfamily GEN};
    \path(11.4,-.3)node(){\strut\ttfamily OCOMP};
    \path(1.4,.3)node(){\strut\large\ttfamily [−]};
    \path(3.4,.3)node(){\strut\large\ttfamily [-]/[+]};
    \path(5.4,.3)node(){\strut\large\ttfamily [-]/[+]};
    \path(7.4,.3)node(){\strut\large\ttfamily -/[+]};
    \path(9.4,.3)node(){\strut\large\ttfamily -/[+]};
    \path(11.4,.3)node(){\strut\large\ttfamily [-]/[+]};
}
\vspace{-6pt} 
\end{mybox}
\vspace{-6pt} 
\textit{Note}. 0 = unrelativizable; − = gapping; + = resumption; [ ] = direct evidence.\par
\setstretch{1.5}\vspace{-12pt}
\end{figure}

Step 0 shows the L2 initial state, on the assumption that the learner has encountered positive evidence in Mandarin for gapping and/or resumption at each level on the scale. When the L2er encounters positive evidence of English gapped RCs at SU in Step~1, no change is needed because the interlanguage already has mandatory gapping at that position. However, restructuring is required when positive evidence of English gapping at IO arrives in Step~2, and the L2er responds by adding gapped RCs at that position. Finally, the L2er responds to positive evidence of English gapping at OCOMP in Step~3 by adding gapping there, with Continuity demanding that gapping also be added at GEN and OBL. The end result is a grammar state that is intermediate between the L2 initial state (from Mandarin) and the target state (English), with mandatory gapping at SU and optionality between gapping and resumption at DO through OCOMP.

The L2er can get to the target grammar state if instead we posit that transferred values lose their `supported by direct evidence' designation, as illustrated in Figure~\ref{fig:overwriting}. I call this the `overwriting' approach because it allows values transferred from the L1 to be overwritten by new values from target-language input.

\begin{figure}
\caption{The overwriting approach to restructuring for L1-Mandarin L2ers of English in single-clause RCs}
\label{fig:overwriting}
\setstretch{1}
\raggedright
\begin{mybox}
{\bfseries Step 0: L2 initial state with full transfer}
\par\vspace{2pt}
Assumed initial state has gapping at SU, optionality at DO, and resumption at IO through OCOMP
\par\vspace{6pt}
\tikz{
    \draw[draw=white, line width=1pt] (0,-5pt) -- (0,5pt);
    \draw[line width=1pt] (0.4,0) -- (12.4,0);
    \draw[line width=1pt] (0.4,-5pt) -- (0.4,5pt);
    \draw[line width=1pt] (2.4,-5pt) -- (2.4,5pt);
    \draw[line width=1pt] (4.4,-5pt) -- (4.4,5pt);
    \draw[line width=1pt] (6.4,-5pt) -- (6.4,5pt);
    \draw[line width=1pt] (8.4,-5pt) -- (8.4,5pt);
    \draw[line width=1pt] (10.4,-5pt) -- (10.4,5pt);
    \draw[line width=1pt] (12.4,-5pt) -- (12.4,5pt);
    \path(1.4,-.3)node(){\strut\ttfamily SU};
    \path(3.4,-.3)node(){\strut\ttfamily DO};
    \path(5.4,-.3)node(){\strut\ttfamily IO};
    \path(7.4,-.3)node(){\strut\ttfamily OBL};
    \path(9.4,-.3)node(){\strut\ttfamily GEN};
    \path(11.4,-.3)node(){\strut\ttfamily OCOMP};
    \path(1.4,.3)node(){\strut\large\ttfamily −};
    \path(3.4,.3)node(){\strut\large\ttfamily -/+};
    \path(5.4,.3)node(){\strut\large\ttfamily +};
    \path(7.4,.3)node(){\strut\large\ttfamily +};
    \path(9.4,.3)node(){\strut\large\ttfamily +};
    \path(11.4,.3)node(){\strut\large\ttfamily +};
}\par
\vspace{3pt}
{\bfseries Step 1: Positive evidence of gapped RCs at SU}
\par\vspace{2pt}
Record positive evidence of gapped RCs at SU
\par\vspace{6pt}
\tikz{
    \draw[draw=white, line width=1pt] (0,-5pt) -- (0,5pt);
    \draw[line width=1pt] (0.4,0) -- (12.4,0);
    \draw[line width=1pt] (0.4,-5pt) -- (0.4,5pt);
    \draw[line width=1pt] (2.4,-5pt) -- (2.4,5pt);
    \draw[line width=1pt] (4.4,-5pt) -- (4.4,5pt);
    \draw[line width=1pt] (6.4,-5pt) -- (6.4,5pt);
    \draw[line width=1pt] (8.4,-5pt) -- (8.4,5pt);
    \draw[line width=1pt] (10.4,-5pt) -- (10.4,5pt);
    \draw[line width=1pt] (12.4,-5pt) -- (12.4,5pt);
    \path(1.4,-.3)node(){\strut\ttfamily SU};
    \path(3.4,-.3)node(){\strut\ttfamily DO};
    \path(5.4,-.3)node(){\strut\ttfamily IO};
    \path(7.4,-.3)node(){\strut\ttfamily OBL};
    \path(9.4,-.3)node(){\strut\ttfamily GEN};
    \path(11.4,-.3)node(){\strut\ttfamily OCOMP};
    \path(1.4,.3)node(){\strut\large\ttfamily [-]};
    \path(3.4,.3)node(){\strut\large\ttfamily -/+};
    \path(5.4,.3)node(){\strut\large\ttfamily +};
    \path(7.4,.3)node(){\strut\large\ttfamily +};
    \path(9.4,.3)node(){\strut\large\ttfamily +};
    \path(11.4,.3)node(){\strut\large\ttfamily +};
}\par\vspace{3pt}
{\bfseries Step 2: Positive evidence of gapped RCs at IO }
\par\vspace{2pt}
Gapping added at IO; Preemption demands that resumption be removed at both IO and DO (because there is direct evidence for gapping at IO and Continuity demands that gapping be available at DO, but nothing forces resumption to be maintained at either of those positions)
\par\vspace{6pt}
\tikz{
    \draw[draw=white, line width=1pt] (0,-5pt) -- (0,5pt);
    \draw[line width=1pt] (0.4,0) -- (12.4,0);
    \draw[line width=1pt] (0.4,-5pt) -- (0.4,5pt);
    \draw[line width=1pt] (2.4,-5pt) -- (2.4,5pt);
    \draw[line width=1pt] (4.4,-5pt) -- (4.4,5pt);
    \draw[line width=1pt] (6.4,-5pt) -- (6.4,5pt);
    \draw[line width=1pt] (8.4,-5pt) -- (8.4,5pt);
    \draw[line width=1pt] (10.4,-5pt) -- (10.4,5pt);
    \draw[line width=1pt] (12.4,-5pt) -- (12.4,5pt);
    \path(1.4,-.3)node(){\strut\ttfamily SU};
    \path(3.4,-.3)node(){\strut\ttfamily DO};
    \path(5.4,-.3)node(){\strut\ttfamily IO};
    \path(7.4,-.3)node(){\strut\ttfamily OBL};
    \path(9.4,-.3)node(){\strut\ttfamily GEN};
    \path(11.4,-.3)node(){\strut\ttfamily OCOMP};
    \path(1.4,.3)node(){\strut\large\ttfamily [-]};
    \path(3.4,.3)node(){\strut\large\ttfamily -};
    \path(5.4,.3)node(){\strut\large\ttfamily [-]};
    \path(7.4,.3)node(){\strut\large\ttfamily +};
    \path(9.4,.3)node(){\strut\large\ttfamily +};
    \path(11.4,.3)node(){\strut\large\ttfamily +};
}\par\vspace{3pt}
{\bfseries Step 3: Positive evidence of gapped RCs at OCOMP}
\par\vspace{2pt}
Gapping added at OCOMP; Preemption demands that resumption be removed at OBL through OCOMP (because there is direct evidence of gapping at OCOMP and Continuity demands that resumption be available at all intermediate positions between SU and OCOMP, but nothing forces resumption to be maintained anywhere in that range)
\par\vspace{6pt}
\tikz{
    \draw[draw=white, line width=1pt] (0,-5pt) -- (0,5pt);
    \draw[line width=1pt] (0.4,0) -- (12.4,0);
    \draw[line width=1pt] (0.4,-5pt) -- (0.4,5pt);
    \draw[line width=1pt] (2.4,-5pt) -- (2.4,5pt);
    \draw[line width=1pt] (4.4,-5pt) -- (4.4,5pt);
    \draw[line width=1pt] (6.4,-5pt) -- (6.4,5pt);
    \draw[line width=1pt] (8.4,-5pt) -- (8.4,5pt);
    \draw[line width=1pt] (10.4,-5pt) -- (10.4,5pt);
    \draw[line width=1pt] (12.4,-5pt) -- (12.4,5pt);
    \path(1.4,-.3)node(){\strut\ttfamily SU};
    \path(3.4,-.3)node(){\strut\ttfamily DO};
    \path(5.4,-.3)node(){\strut\ttfamily IO};
    \path(7.4,-.3)node(){\strut\ttfamily OBL};
    \path(9.4,-.3)node(){\strut\ttfamily GEN};
    \path(11.4,-.3)node(){\strut\ttfamily OCOMP};
    \path(1.4,.3)node(){\strut\large\ttfamily [-]};
    \path(3.4,.3)node(){\strut\large\ttfamily -};
    \path(5.4,.3)node(){\strut\large\ttfamily [-]};
    \path(7.4,.3)node(){\strut\large\ttfamily -};
    \path(9.4,.3)node(){\strut\large\ttfamily -};
    \path(11.4,.3)node(){\strut\large\ttfamily [-]};
}
\vspace{-6pt} 
\end{mybox}
\vspace{-6pt} 
\textit{Note}. 0 = unrelativizable; − = gapping; + = resumption; [ ] = direct evidence.\par
\setstretch{1.5}\vspace{-12pt}
\end{figure}

At Step~1, all the learner needs to do is record that there is now positive evidence of gapped RCs at SU. A more dramatic form of restructuring takes place at Step~2, where the learner encounters positive evidence of gapped RCs at IO. Preemption demands that gapping overwrite resumption at IO because whereas there is direct evidence for gapping there, nothing forces resumption to be maintained in the same position. A similar process takes place at DO: Gapping must be added to satisfy Continuity, but resumption is pushed out because there is nothing forcing it to be maintained there. Finally, the same overwriting process is repeated down to the bottom of the scale in Step~3 when the L2er receives positive evidence of gapping at OCOMP. The end result is the target state with mandatory gapping across the entire scale.

Figure~\ref{fig:overwriting-ko} lays out what happens when we apply the overwriting approach to the case of an L1-English L2er of Korean. Recall that Korean has no relativization at OCOMP, so in this case reaching the target grammar requires not only the overwriting of one relativization option with another but also the complete removal of relativization from one level on the scale.

\begin{figure}
\caption{The overwriting approach to restructuring for L1-English L2ers of Korean in single-clause RCs}
\label{fig:overwriting-ko}
\setstretch{1}
\raggedright
\begin{mybox}
{\bfseries Step 0: L2 initial state with full transfer}
\par\vspace{2pt}
Assumed initial state has gapping across the entire scale
\par\vspace{6pt}
\tikz{
    \draw[draw=white, line width=1pt] (0,-5pt) -- (0,5pt);
    \draw[line width=1pt] (0.4,0) -- (12.4,0);
    \draw[line width=1pt] (0.4,-5pt) -- (0.4,5pt);
    \draw[line width=1pt] (2.4,-5pt) -- (2.4,5pt);
    \draw[line width=1pt] (4.4,-5pt) -- (4.4,5pt);
    \draw[line width=1pt] (6.4,-5pt) -- (6.4,5pt);
    \draw[line width=1pt] (8.4,-5pt) -- (8.4,5pt);
    \draw[line width=1pt] (10.4,-5pt) -- (10.4,5pt);
    \draw[line width=1pt] (12.4,-5pt) -- (12.4,5pt);
    \path(1.4,-.3)node(){\strut\ttfamily SU};
    \path(3.4,-.3)node(){\strut\ttfamily DO};
    \path(5.4,-.3)node(){\strut\ttfamily IO};
    \path(7.4,-.3)node(){\strut\ttfamily OBL};
    \path(9.4,-.3)node(){\strut\ttfamily GEN};
    \path(11.4,-.3)node(){\strut\ttfamily OCOMP};
    \path(1.4,.3)node(){\strut\large\ttfamily -};
    \path(3.4,.3)node(){\strut\large\ttfamily -};
    \path(5.4,.3)node(){\strut\large\ttfamily -};
    \path(7.4,.3)node(){\strut\large\ttfamily -};
    \path(9.4,.3)node(){\strut\large\ttfamily -};
    \path(11.4,.3)node(){\strut\large\ttfamily -};
}\par
\vspace{3pt}
{\bfseries Step 1: Positive evidence of gapped RCs at SU}
\par\vspace{2pt}
Record positive evidence of gapped RCs at SU
\par\vspace{6pt}
\tikz{
    \draw[draw=white, line width=1pt] (0,-5pt) -- (0,5pt);
    \draw[line width=1pt] (0.4,0) -- (12.4,0);
    \draw[line width=1pt] (0.4,-5pt) -- (0.4,5pt);
    \draw[line width=1pt] (2.4,-5pt) -- (2.4,5pt);
    \draw[line width=1pt] (4.4,-5pt) -- (4.4,5pt);
    \draw[line width=1pt] (6.4,-5pt) -- (6.4,5pt);
    \draw[line width=1pt] (8.4,-5pt) -- (8.4,5pt);
    \draw[line width=1pt] (10.4,-5pt) -- (10.4,5pt);
    \draw[line width=1pt] (12.4,-5pt) -- (12.4,5pt);
    \path(1.4,-.3)node(){\strut\ttfamily SU};
    \path(3.4,-.3)node(){\strut\ttfamily DO};
    \path(5.4,-.3)node(){\strut\ttfamily IO};
    \path(7.4,-.3)node(){\strut\ttfamily OBL};
    \path(9.4,-.3)node(){\strut\ttfamily GEN};
    \path(11.4,-.3)node(){\strut\ttfamily OCOMP};
    \path(1.4,.3)node(){\strut\large\ttfamily [-]};
    \path(3.4,.3)node(){\strut\large\ttfamily -};
    \path(5.4,.3)node(){\strut\large\ttfamily -};
    \path(7.4,.3)node(){\strut\large\ttfamily -};
    \path(9.4,.3)node(){\strut\large\ttfamily -};
    \path(11.4,.3)node(){\strut\large\ttfamily -};
}\par\vspace{3pt}
{\bfseries Step 2: Positive evidence of gapped RCs at IO}
\par\vspace{2pt}
Record positive evidence of gapped RCs at IO
\par\vspace{6pt}
\tikz{
    \draw[draw=white, line width=1pt] (0,-5pt) -- (0,5pt);
    \draw[line width=1pt] (0.4,0) -- (12.4,0);
    \draw[line width=1pt] (0.4,-5pt) -- (0.4,5pt);
    \draw[line width=1pt] (2.4,-5pt) -- (2.4,5pt);
    \draw[line width=1pt] (4.4,-5pt) -- (4.4,5pt);
    \draw[line width=1pt] (6.4,-5pt) -- (6.4,5pt);
    \draw[line width=1pt] (8.4,-5pt) -- (8.4,5pt);
    \draw[line width=1pt] (10.4,-5pt) -- (10.4,5pt);
    \draw[line width=1pt] (12.4,-5pt) -- (12.4,5pt);
    \path(1.4,-.3)node(){\strut\ttfamily SU};
    \path(3.4,-.3)node(){\strut\ttfamily DO};
    \path(5.4,-.3)node(){\strut\ttfamily IO};
    \path(7.4,-.3)node(){\strut\ttfamily OBL};
    \path(9.4,-.3)node(){\strut\ttfamily GEN};
    \path(11.4,-.3)node(){\strut\ttfamily OCOMP};
    \path(1.4,.3)node(){\strut\large\ttfamily [-]};
    \path(3.4,.3)node(){\strut\large\ttfamily -};
    \path(5.4,.3)node(){\strut\large\ttfamily [-]};
    \path(7.4,.3)node(){\strut\large\ttfamily -};
    \path(9.4,.3)node(){\strut\large\ttfamily -};
    \path(11.4,.3)node(){\strut\large\ttfamily -};
}\par\vspace{3pt}
{\bfseries Step 3: Positive evidence of resumptive RCs at GEN}
\par\vspace{2pt}
Resumption added at GEN; Preemption demands that only resumption be available at GEN (because nothing forces gapping to be maintained there); Continuity forbids gapping from being maintained at OCOMP, so gapping is removed from that position.
\par\vspace{6pt}
\tikz{
    \draw[draw=white, line width=1pt] (0,-5pt) -- (0,5pt);
    \draw[line width=1pt] (0.4,0) -- (12.4,0);
    \draw[line width=1pt] (0.4,-5pt) -- (0.4,5pt);
    \draw[line width=1pt] (2.4,-5pt) -- (2.4,5pt);
    \draw[line width=1pt] (4.4,-5pt) -- (4.4,5pt);
    \draw[line width=1pt] (6.4,-5pt) -- (6.4,5pt);
    \draw[line width=1pt] (8.4,-5pt) -- (8.4,5pt);
    \draw[line width=1pt] (10.4,-5pt) -- (10.4,5pt);
    \draw[line width=1pt] (12.4,-5pt) -- (12.4,5pt);
    \path(1.4,-.3)node(){\strut\ttfamily SU};
    \path(3.4,-.3)node(){\strut\ttfamily DO};
    \path(5.4,-.3)node(){\strut\ttfamily IO};
    \path(7.4,-.3)node(){\strut\ttfamily OBL};
    \path(9.4,-.3)node(){\strut\ttfamily GEN};
    \path(11.4,-.3)node(){\strut\ttfamily OCOMP};
    \path(1.4,.3)node(){\strut\large\ttfamily [-]};
    \path(3.4,.3)node(){\strut\large\ttfamily -};
    \path(5.4,.3)node(){\strut\large\ttfamily [-]};
    \path(7.4,.3)node(){\strut\large\ttfamily -};
    \path(9.4,.3)node(){\strut\large\ttfamily [+]};
    \path(11.4,.3)node(){\strut\ttfamily 0};
}
\vspace{-6pt} 
\end{mybox}
\vspace{-6pt} 
\textit{Note}. 0 = unrelativizable; − = gapping; + = resumption; [ ] = direct evidence.\par
\setstretch{1.5}\vspace{-12pt}
\end{figure}

Step~0 shows the L2 initial state assuming full transfer from English, with mandatory gapping across the entire scale. In Steps~1 and 2, the learner simply records that there is now direct evidence of gapped RCs at SU and OBL. In Step~3, the L2er responds to positive evidence of resumptive RCs at GEN by overwriting gapping with resumption there. Note that Affordability does not require that resumption be extended upwards from GEN since another option (i.e.,~gapping) is already available at all higher positions on the scale. The learner is then forced to remove gapping from OCOMP because Continuity forbids an option from operating on disjointed segments of the scale. Also, resumption cannot take the place of gapping at OCOMP because extending resumption downward from GEN without positive evidence would violate Conservatism. The result is therefore the target grammar state with mandatory gapping at SU through OBL, mandatory resumption at GEN, and no relativization at OCOMP.

Now let us consider what happens when we apply the additive and overwriting approaches to the SLW hierarchy for ORCs. Figure~\ref{fig:slw-additive} shows how the acquisition process might unfold according to the additive approach for an L1-Mandarin L2er of English. 

\begin{figure}
\caption{The additive approach for L1-Mandarin L2ers on the SLW hierarchy for English ORCs}
\label{fig:slw-additive}
\setstretch{1}
\raggedright
\begin{mybox}
{\bfseries Step 0: L2 initial state with full transfer}
\par\vspace{2pt}
Assumed initial state has optionality between gapping and resumption at S through L and no relativization at W
\par\vspace{6pt}
\tikz{
    \draw[draw=white, line width=1pt] (0,-5pt) -- (0,5pt);
    \draw[line width=1pt] (0.4,0) -- (12.4,0);
    \draw[line width=1pt] (0.4,-5pt) -- (0.4,5pt);
    \draw[line width=1pt] (4.4,-5pt) -- (4.4,5pt);
    \draw[line width=1pt] (8.4,-5pt) -- (8.4,5pt);
    \draw[line width=1pt] (12.4,-5pt) -- (12.4,5pt);
    \path(2.4,-.3)node(){\strut\ttfamily\strut S};
    \path(6.4,-.3)node(){\strut\ttfamily\strut L};
    \path(10.4,-.3)node(){\strut\ttfamily\strut W};
    \path(2.4,.3)node(){\strut\large\ttfamily [-]/[+]};
    \path(6.4,.3)node(){\strut\large\ttfamily [-]/[+]};
    \path(10.4,.3)node(){\strut\ttfamily 0};
}\par
\vspace{3pt}
{\bfseries Step 1: Positive evidence of gapping in short-distance ORCs}
\par\vspace{2pt}
No change
\par\vspace{6pt}
\tikz{
    \draw[draw=white, line width=1pt] (0,-5pt) -- (0,5pt);
    \draw[line width=1pt] (0.4,0) -- (12.4,0);
    \draw[line width=1pt] (0.4,-5pt) -- (0.4,5pt);
    \draw[line width=1pt] (4.4,-5pt) -- (4.4,5pt);
    \draw[line width=1pt] (8.4,-5pt) -- (8.4,5pt);
    \draw[line width=1pt] (12.4,-5pt) -- (12.4,5pt);
    \path(2.4,-.3)node(){\strut\ttfamily\strut S};
    \path(6.4,-.3)node(){\strut\ttfamily\strut L};
    \path(10.4,-.3)node(){\strut\ttfamily\strut W};
    \path(2.4,.3)node(){\strut\large\ttfamily [-]/[+]};
    \path(6.4,.3)node(){\strut\large\ttfamily [-]/[+]};
    \path(10.4,.3)node(){{\ttfamily 0}};
}\par
\vspace{3pt}
{\bfseries Step 2: Positive evidence of gapping in long-distance ORCs}
\par\vspace{2pt}
No change
\par\vspace{6pt}
\tikz{
    \draw[draw=white, line width=1pt] (0,-5pt) -- (0,5pt);
    \draw[line width=1pt] (0.4,0) -- (12.4,0);
    \draw[line width=1pt] (0.4,-5pt) -- (0.4,5pt);
    \draw[line width=1pt] (4.4,-5pt) -- (4.4,5pt);
    \draw[line width=1pt] (8.4,-5pt) -- (8.4,5pt);
    \draw[line width=1pt] (12.4,-5pt) -- (12.4,5pt);
    \path(2.4,-.3)node(){\strut\ttfamily\strut S};
    \path(6.4,-.3)node(){\strut\ttfamily\strut L};
    \path(10.4,-.3)node(){\strut\ttfamily\strut W};
    \path(2.4,.3)node(){\strut\large\ttfamily [-]/[+]};
    \path(6.4,.3)node(){\strut\large\ttfamily [-]/[+]};
    \path(10.4,.3)node(){\strut\ttfamily 0};
}
\vspace{-6pt} 
\end{mybox}
\vspace{-6pt} 
\textit{Note}. 0 = unrelativizable; − = gapping; + = resumption; [ ] = direct evidence.\par
\setstretch{1.5}\vspace{-12pt}
\end{figure}

\noindent Step~0 shows the L2 initial state under full transfer. Because extending gapped or resumptive RCs to a newly-relevant level at the bottom of the scale would violate Conservatism, I assume here that the `W' level on the scale, which is N/A~(`not applicable') in Mandarin, defaults to 0~(`no relativization') at the L2 initial state. No restructuring takes place when positive evidence of gapping at the short-distance and long-distance levels arrives in Steps~1 and 2, respectively, because the grammar already has gapping there. Meanwhile, resumption persists as an option for relativization in those environments because the `supported by direct evidence' designation is carried over from the L1 grammar. Positive evidence of gapping at the \textit{wh}-island level never arrives—at least not in the quantity needed to trigger restructuring---and so the \textit{wh}-island level on the scale remains set to `no relativization,' just as it is for L1-English adults. The L2er therefore ends up with a nontarget grammar state with optionality between gapping and resumption at the short-distance and long-distance levels and no relativization at the \textit{wh}-island level.

Things play out very differently with the overwriting approach to restructuring, as demonstrated in Figure~\ref{fig:slw-overwriting}. 

\begin{figure}
\caption{The overwriting approach for L1-Mandarin L2ers on the SLW hierarchy for English ORCs}
\label{fig:slw-overwriting}
\setstretch{1}
\raggedright
\begin{mybox}
{\bfseries Step 0: L2 initial state with full transfer}
\par\vspace{2pt}
Assumed initial state with optionality between gapping and resumption at S through L and no relativization at W
\par\vspace{6pt}
\tikz{
    \draw[draw=white, line width=1pt] (0,-5pt) -- (0,5pt);
    \draw[line width=1pt] (0.4,0) -- (12.4,0);
    \draw[line width=1pt] (0.4,-5pt) -- (0.4,5pt);
    \draw[line width=1pt] (4.4,-5pt) -- (4.4,5pt);
    \draw[line width=1pt] (8.4,-5pt) -- (8.4,5pt);
    \draw[line width=1pt] (12.4,-5pt) -- (12.4,5pt);
    \path(2.4,-.3)node(){\strut\ttfamily\strut S};
    \path(6.4,-.3)node(){\strut\ttfamily\strut L};
    \path(10.4,-.3)node(){\strut\ttfamily\strut W};
    \path(2.4,.3)node(){\strut\large\ttfamily -/+};
    \path(6.4,.3)node(){\strut\large\ttfamily -/+};
    \path(10.4,.3)node(){\strut\ttfamily 0};
}\par
\vspace{3pt}
{\bfseries Step 1: Positive evidence of gapping in short-distance ORCs}
\par\vspace{2pt}
Positive evidence of gapping recorded at S; Preemption demands that only gapping be available at S (because nothing forces resumption to be maintained there) 
\par\vspace{6pt}
\tikz{
    \draw[draw=white, line width=1pt] (0,-5pt) -- (0,5pt);
    \draw[line width=1pt] (0.4,0) -- (12.4,0);
    \draw[line width=1pt] (0.4,-5pt) -- (0.4,5pt);
    \draw[line width=1pt] (4.4,-5pt) -- (4.4,5pt);
    \draw[line width=1pt] (8.4,-5pt) -- (8.4,5pt);
    \draw[line width=1pt] (12.4,-5pt) -- (12.4,5pt);
    \path(2.4,-.3)node(){\strut\ttfamily\strut S};
    \path(6.4,-.3)node(){\strut\ttfamily\strut L};
    \path(10.4,-.3)node(){\strut\ttfamily\strut W};
    \path(2.4,.3)node(){\strut\large\ttfamily [-]};
    \path(6.4,.3)node(){\strut\large\ttfamily -/+};
    \path(10.4,.3)node(){{\ttfamily 0}};
}\par
\vspace{3pt}
{\bfseries Step 2: Positive evidence of gapping in long-distance ORCs}
\par\vspace{2pt}
Positive evidence of gapping recorded at L; Preemption demands that only gapping be available at L (because nothing forces resumption to be maintained there) 
\par\vspace{6pt}
\tikz{
    \draw[draw=white, line width=1pt] (0,-5pt) -- (0,5pt);
    \draw[line width=1pt] (0.4,0) -- (12.4,0);
    \draw[line width=1pt] (0.4,-5pt) -- (0.4,5pt);
    \draw[line width=1pt] (4.4,-5pt) -- (4.4,5pt);
    \draw[line width=1pt] (8.4,-5pt) -- (8.4,5pt);
    \draw[line width=1pt] (12.4,-5pt) -- (12.4,5pt);
    \path(2.4,-.3)node(){\strut\ttfamily\strut S};
    \path(6.4,-.3)node(){\strut\ttfamily\strut L};
    \path(10.4,-.3)node(){\strut\ttfamily\strut W};
    \path(2.4,.3)node(){\strut\large\ttfamily [-]};
    \path(6.4,.3)node(){\strut\large\ttfamily [-]};
    \path(10.4,.3)node(){\strut\ttfamily 0};
}
\vspace{-6pt} 
\end{mybox}
\vspace{-6pt} 
\textit{Note}. 0 = unrelativizable; − = gapping; + = resumption; [ ] = direct evidence.\par
\setstretch{1.5}\vspace{-12pt}
\end{figure}

\noindent At Step~1, positive evidence for gapped RCs at the short-distance level does not conflict with the initial L2 grammar state, which has optionality between gapping and resumption in that environment; however, recall that Preemption only permits optionality where it is forced. Since neither positive evidence nor Continuity forces resumption to be maintained at S, Preemption demands that it be removed from that environment. The same overwriting process occurs at Step 2, where the L2er responds to positive evidence of gapped RCs at the long-distance level by purging resumption from that environment. The end result is the target state with mandatory gapping at S through L and no relativization at W. Note that this approach does not rely on indirect negative evidence (as conceptualized by, e.g.,~\cite{plough1994}), but rather on a Preemption principle which tolerates optionality only when forced to do so (by positive evidence or by Continuity). 

In examining judgment data from L1-Korean vs. L1-Mandarin L2ers of English at a range of proficiency levels, the present dissertation offers a chance to test the predictions of the additive and overwriting approaches. If at least some advanced L1-Mandarin L2ers consistently reject resumption in one or more environments where it is permitted in the L1 (short-distance ORCs and long-distance ORCs/SRCs), this would suggest that they have purged resumption from those environments, as would be expected under the overwriting approach to restructuring. It seems reasonable to expect that levels of attainment in this regard might be higher in the short-distance environment (where relevant input is assumed to be relatively abundant) than in the long-distance environment (where relevant input is assumed to be somewhat less frequent). If, however, L1-Mandarin L2ers continue to accept resumption in domains where it is permitted in the L1 even at high levels of English proficiency, this would be consistent with the additive approach to restructuring. Furthermore, if at least some L2ers consistently accept resumption in environments where it is not permitted in either the L1 or the target language (short-distance SRCs for L1-Mandarin L2ers and both short- and long-distance ORCs/SRCs for L1-Korean L2ers), this would suggest that resumption can enter English interlanguage grammars even in the absence of L1 influence, in line with the proposals made by \textcite{hyltenstam1984,hyltenstam1990} and \textcite{eckman2004}.

%========================================================%
\section{Summing up}\label{acq-conclusion}
%========================================================%

In this chapter, I provided an overview of some of the most influential studies that have touched on the L2 acceptability and processing of gapped RCs vs. resumptive RCs. I then proceeded to outline an approach to the L1 and L2 development of implicit knowledge of the distribution of these RC types in which five principles (Gatekeeping, Conservatism, Affordability, Continuity, and Preemption) operate in conjunction with hierarchies of processing difficulty and positive evidence in the input to guide learners toward constructing the target grammar state. I also offered two distinct approaches to how changes in the L2 grammar might be effectuated: an additive approach (which does not allow for relativization options transferred from the L1 to be purged from the interlanguage) and an overwriting approach (which does allow such purging to occur).

The chapter also introduced several key objectives of this dissertation. First, I aim to supplement existing research on L2 resumption in RCs by examining an understudied set of processing asymmetries for RCs, which I call the short-distance vs. long-distance vs. \textit{wh}-island (SLW) hierarchy. Second, I test the Performance-Competence Correspondence Assumption (PCCA) by investigating whether L2ers' resumptive RCs in English are best categorized as grammatical resumptives or processing resumptives in the interlanguages of L1-Korean L2ers and L1-Mandarin L2ers. In the event that some L2ers do consistently accept resumptive RCs in English, I also investigate whether that RP acceptance is generally traceable to L1 transfer or whether there might be an alternative source for it. Finally, I test whether such L2ers can become sensitive to the English prohibition on resumption in RCs over time.

This concludes the background chapters of the dissertation. In the next chapter, I shift the focus to the design of my experiments.